\documentclass[11pt,a4paper]{article}

% Packages
\usepackage[utf8]{inputenc}
\usepackage[T1]{fontenc}
\usepackage{amsmath,amssymb,amsfonts}
\usepackage{graphicx}
\usepackage{booktabs}
\usepackage{array}
\usepackage{multirow}
\usepackage{geometry}
\usepackage{setspace}
\usepackage{gensymb}
\usepackage{siunitx}
\usepackage{textcomp}
\usepackage{textgreek}
\usepackage{threeparttable}
\usepackage{hyperref}
\usepackage{xcolor}
\usepackage{float}
\usepackage{caption}
\usepackage{subcaption}
\usepackage{tcolorbox}
\usepackage{longtable}

% Unicode character support (preserved from original)
\usepackage{newunicodechar}
\newunicodechar{⁵}{\textsuperscript{5}}
\newunicodechar{⁴}{\textsuperscript{4}}
\newunicodechar{⁰}{\textsuperscript{0}}
\newunicodechar{ⁿ}{\textsuperscript{n}}
\newunicodechar{≤}{\leq}
\newunicodechar{≈}{\approx}
\newunicodechar{−}{\ensuremath{-}}
\newunicodechar{₀}{$_0$}
\newunicodechar{₁}{$_1$}
\newunicodechar{₂}{$_2$}
\newunicodechar{₃}{$_3$}
\newunicodechar{₄}{$_4$}
\newunicodechar{₅}{$_5$}
\newunicodechar{₆}{$_6$}
\newunicodechar{₇}{$_7$}
\newunicodechar{₈}{$_8$}
\newunicodechar{₉}{$_9$}
\newunicodechar{⁷}{\textsuperscript{7}}
\newunicodechar{√}{\sqrt{}}
\newunicodechar{ⱼ}{\textsubscript{j}}
\newunicodechar{ₖ}{_k}
\newunicodechar{≥}{\geq{}}
\newunicodechar{⁶}{\textsuperscript{6}}
\newunicodechar{⁻}{\textsuperscript{-}}

% Language support
\usepackage[russian,english]{babel}

% Page setup
\geometry{margin=1in}
\onehalfspacing

% Hyperref setup
\hypersetup{
    colorlinks=true,
    linkcolor=blue,
    citecolor=blue,
    urlcolor=blue
}

% Custom commands
\newcommand{\phisym}{\varphi}
\newcommand{\fzero}{f_0}

% Colored boxes for key equations
\newtcolorbox{keyequation}{
    colback=blue!5!white,
    colframe=blue!75!black,
    fonttitle=\bfseries,
    title=Key Equation
}

\newtcolorbox{keyinsight}{
    colback=green!5!white,
    colframe=green!50!black,
    fonttitle=\bfseries,
    title=Key Insight
}

\newtcolorbox{criticalpoint}{
    colback=red!5!white,
    colframe=red!50!black,
    fonttitle=\bfseries,
    title=Critical Methodological Point
}

% Bibliography (biblatex - preserved from original)
\usepackage[style=apa,backend=biber]{biblatex}
\addbibresource{references.bib}
\addbibresource{intro.bib}
\addbibresource{methods.bib}
\addbibresource{results.bib}
\addbibresource{discussion.bib}

\begin{document}

% Title
\begin{center}
{\LARGE \textbf{Transient High-Coherence Neural States with Golden Ratio Harmonic Organization at Earth-Resonant Frequencies}}\\[0.5cm]
{\large \textbf{Discovery of Schumann Ignition Events in Human EEG}}\\[1cm]

{\large Michael Lacy}\\[0.5cm]
{\small michael.lacy@gmail.com}\\[0.5cm]

\today
\end{center}

\vspace{0.5cm}

We report the discovery and characterization of Schumann Ignition Events (SIEs)—transient, high-coherence neural states occurring at frequencies within the range of Earth's Schumann Resonance. Analysis of 1,366 events across 91 participants, 661 recording sessions, and three consumer EEG devices (Emotiv EPOC X, Muse, Insight) spanning five cognitive contexts reveals brief episodes of multi-band network synchronization with consistent spectral architecture.

SIEs are characterized by elevated power and phase coherence across multiple frequency bands simultaneously. The fundamental frequency (7.63 ± 0.33 Hz) falls within the natural fluctuation range of the first Schumann Resonance mode (7.0–8.2 Hz). At this frequency, 81.5\% of events exhibited phase-locking values exceeding 0.6, and 60.5\% showed magnitude-squared coherence above 0.5. Events followed a stereotyped six-phase temporal evolution in which coherence elevation preceded amplitude increases—a signature consistent with network-level coordination rather than local oscillatory bursts. Detected frequencies converged across devices and cognitive contexts for lower harmonics (SR1, SR3), supporting measurement validity, though higher harmonics showed modest device-related variability.

Harmonic frequency analysis revealed that ratios between detected harmonics approximate powers of the golden ratio (φ = 1.618) with less than 1\% error. Critically, individual harmonic frequencies varied independently across events (all pairwise |r| < 0.04), yet ratio precision was preserved—indicating that golden ratio relationships are embedded in population-level frequency distributions rather than arising from within-event coupling. Null controls confirmed genuine phase synchronization (phase-randomized surrogates: d = 2.20) and spectral exceptionality versus random triplets from actual EEG peaks (p $<$ 0.0001, d = 1.44). Baseline periods outside detected events showed identical frequency ratios with reduced power and coherence, establishing that SIEs represent transient amplification of continuously present oscillatory structure.

These findings establish SIEs as a reproducible multi-band synchronization phenomenon with precise frequency organization corresponding to the Schumann Resonance range. Whether the SR-neural frequency correspondence reflects functional coupling, convergent biophysical constraints, or coincidence cannot be determined without concurrent geomagnetic measurement. The golden ratio harmonic architecture—present continuously but amplified during SIEs—suggests intrinsic organizational principles in neural oscillatory dynamics warranting further investigation.
\newpage
\tableofcontents
\newpage
\section{Introduction}
Neural oscillations are fundamental to brain function, coordinating activity across distributed networks to support perception, cognition, and behavior ( \cite{gy_rgy_buzs_ki_c98dac5c}; \cite{pascal_fries_7b9d1ced}). These rhythmic fluctuations in neuronal population activity occur across characteristic frequency bands—from slow delta oscillations (<4 Hz) through high gamma (>80 Hz)—each associated with distinct functional roles (\cite{wolfgang_klimesch_be1f7786}; \cite{engel2010beta}, \cite{cohen2017where}). Although substantial research has characterized sustained oscillatory states and their behavioral correlates, transient synchronization phenomena—brief episodes of enhanced coordination across brain regions—have received comparatively less systematic investigation despite their potential significance for understanding rapid neural state transitions and information integration (\cite{siegel2012spectral}; \cite{canolty2010functional}).

Well-characterized examples of transient oscillatory events include sleep spindles (\cite{fernandez2020sleep}), hippocampal sharp-wave ripples (\cite{buzsaki2015hippocampal}), and task-related gamma bursts (\cite{fries2001modulation}). These phenomena share common features: discrete temporal boundaries, elevated power relative to baseline, and enhanced inter-regional synchronization. Such transient events may carry functional significance distinct from sustained oscillations, potentially serving as temporal windows for memory consolidation, inter-regional communication, or state transitions (\cite{john_lisman_100e4cf6}).

The Earth's ionospheric cavity supports a set of global electromagnetic resonances, termed Schumann Resonances (SR), arising from electromagnetic waves circling the Earth-ionosphere waveguide driven primarily by global lightning activity (\cite{schumann1952strahlungslosen}; \cite{nickolaenko2014schumann}). The fundamental SR mode occurs at approximately 7.83 Hz, with higher harmonics at approximately 14.3, 20.8, 27.3, and 33.8 Hz—though these values fluctuate daily between approximately 7.0 and 8.2 Hz for the fundamental, depending on ionospheric conditions and global lightning distribution (\cite{price2016global}). The overlap between SR frequencies and canonical EEG bands—particularly theta (4–8 Hz), alpha (8–13 Hz), and beta (13–30 Hz)—has prompted investigation into potential brain-environment relationships, beginning with König's early observations of spectral similarities between SR and human EEG (\cite{konig1974eeg}) and continuing through more recent quantitative analyses (\cite{cherry2002schumann}; \cite{persinger2008possible}; \cite{saroka2016similar}).

Several studies have reported correlations between geomagnetic activity at SR frequencies and various physiological and behavioral measures, including heart rate variability, reaction times, and blood pressure dynamics (\cite{pobachenko2006contingency}; \cite{babayev2007effects}). However, despite this foundational work, critical methodological gaps limit our understanding of SR-aligned neural phenomena. First, prior studies have mainly employed correlational designs examining continuous SR-EEG relationships rather than systematically detecting and characterizing discrete events; the temporal structure, duration, and network dynamics of individual SR-frequency episodes remain poorly characterized. Second, methodological concerns persist: many studies used single EEG devices without cross-platform validation, few employed rigorous null controls to distinguish genuine synchronization from spectral coincidence, and the potential circularity of searching for predefined SR frequencies has not been adequately addressed. Third, harmonic relationships between SR-aligned neural frequencies have received limited quantitative analysis—particularly whether observed frequency ratios reflect electromagnetic cavity physics, alternative mathematical constraints, or measurement variability.

The present study addresses these gaps through systematic detection and characterization of transient high-coherence neural events occurring at frequencies within the Schumann Resonance range. During exploratory analysis of EEG recordings from meditation sessions collected between 2019 and 2024, we repeatedly observed transient periods of elevated power spectral density at approximately 7.6–7.8 Hz across multiple brain regions. The correspondence with the SR frequency range prompted systematic investigation. We discovered that these events involve coordinated power increases at multiple frequencies, accompanied by marked elevation in inter-channel coherence and phase synchronization—suggesting network-wide coordination rather than isolated oscillatory bursts.

We developed a multi-stage detection pipeline and validated the phenomenon across two independent datasets: longitudinal meditation recordings and a publicly available cognitive task dataset (\cite{albuquerque2024physf}). To establish event validity, we implemented multiple null control analyses including phase-randomized surrogates, random temporal window sampling, random frequency triplet comparisons, and shuffled bootstrap tests.

Our analyses reveal 1,366 discrete events—which we term Schumann Ignition Events (SIEs) —exhibiting consistent spectral, temporal, and network characteristics across 91 subjects, three EEG devices, and five cognitive contexts. Events show robust network synchronization (phase-locking values exceeding 0.6 in 81.5\% of events) with coherence elevation preceding power amplification—a temporal signature suggesting network-level coordination. Harmonic frequency ratios approximate powers of the golden ratio (φ = 1.618) with less than 1\% error, a precision validated as spectrally exceptional compared to random triplets sampled from actual EEG spectral peaks (Cohen's d = 1.44, p $<$ 0.0001). Critically, null control analysis revealed that SR-frequency organization exists continuously at baseline with maintained frequency ratios but degraded power and coherence; SIEs represent transient amplification of this persistent oscillatory substrate rather than de novo generation of SR-aligned activity.

These findings establish SIEs as a reproducible neural phenomenon with consistent frequency organization across subjects and recording conditions. The frequencies detected fall within the natural fluctuation range of Earth's Schumann Resonances (\~7.0–8.2 Hz), though whether this correspondence reflects functional brain-environment coupling, convergent biophysical constraints, or coincidence cannot be determined without concurrent geomagnetic measurement. We discuss implications for understanding transient synchronization phenomena and the organizational principles underlying neural frequency architecture.

\section{Methods}

\subsection{Overview}

This study employed a two-phase design: (1) discovery and initial characterization of Schumann Ignition Events (SIEs) in longitudinal meditation EEG recordings, and (2) validation and extension in an independent dataset collected during cognitive tasks. All analyses were performed in Python 3.9+ using NumPy, SciPy, MNE-Python (\cite{harris2020numpy}; \cite{virtanen2020scipy}; \cite{gramfort2013mne}), and the FOOOF library (\cite{donoghue2020fooof}). Complete analysis code and parameter specifications will be made available upon publication.
\subsection{Ethics Statement}
\subsubsection{Discovery Dataset Ethics}
The longitudinal meditation recordings were collected by the first author from their own personal meditation practice for the purpose of self-directed neurofeedback exploration. As self-experimentation by the researcher on their own data, this component does not constitute human subjects research requiring Institutional Review Board (IRB) approval under 45 CFR 46.102. No identifying information beyond that necessary for scientific reporting is included.

\subsubsection{AI Tool Disclosure}
In accordance with policy on AI-assisted tools, we disclose the following: AI tools (Claude Opus 4.5, Anthropic) were used for coding assistance in developing the analysis pipeline and for manuscript preparation. Overleaf was used for paper creation. ChatGPT (OpenAI), Google Gemini, and Grok (xAI) were used as manuscript reviewers. All scientific hypotheses, experimental design, data collection, interpretation of results, and intellectual contributions are the authors' own work. AI-generated code was reviewed, tested, and validated by the author prior to use.

\subsection{Participants and Recordings}

\subsubsection{Discovery Dataset: Longitudinal Meditation Recordings}
The initial discovery and characterization of SIEs was conducted using EEG recordings from a single experienced meditation practitioner (male, age range 44–48 years, >10 years meditation experience). Recordings were collected during eyes-closed meditation sessions spanning November 2019 to March 2024, comprising 34 sessions with a total recording duration of approximately 7.5 hours. Session durations ranged from 2 to 30 minutes (mean ± SD: 13.2 ± 8.7 minutes). Three consumer-grade EEG devices were used across sessions to enable cross-device validation:

\begin{itemize}
    \item \textbf{Muse (2016 model; InteraXon Inc., Toronto, Canada)}- Channels: 4 (TP9, AF7, AF8, TP10; dry electrodes)- Sampling rate: 256 Hz- Reference: FPz- Hardware filters: 0.1–55 Hz bandpass- Sessions recorded: 17

    \item \textbf{Emotiv EPOC X (Emotiv Inc., San Francisco, USA)}- Channels: 14 (AF3, F7, F3, FC5, T7, P7, O1, O2, P8, T8, FC6, F4, F8, AF4; saline-soaked felt pads)- Sampling rate: 128 Hz- Reference: CMS/DRL (P3/P4 locations)- Hardware filters: 0.16–43 Hz bandpass, 50/60 Hz notch- Sessions recorded: 5

    \item \textbf{Emotiv Insight (Emotiv Inc.)}- Channels: 5 (AF3, AF4, T7, T8, Pz; semi-dry polymer sensors)- Sampling rate: 128 Hz- Reference: CMS/DRL- Hardware filters: 0.16–43 Hz bandpass, 50/60 Hz notch- Sessions recorded: 12
\end{itemize}

All recordings were collected in indoor residential environments without electromagnetic shielding. Participants sat in a comfortable upright position with eyes closed. No specific meditation instructions were given beyond maintaining relaxed alertness.

\subsubsection{Validation Dataset 1: PhySF Cognitive Flow Study}
To test whether SIEs generalize beyond meditation to other cognitive states, we analyzed the publicly available Physiological States of Flow (\cite{albuquerque2024physf}). This dataset contains multimodal physiological recordings from participants performing mathematical tasks designed to induce flow states (optimal experience characterized by absorbed attention and intrinsic motivation) along with non-flow control conditions.

\textbf{Original Study Design (Albuquerque et al., 2024)-} N = 26 participants; Task: Arithmetic problem-solving; Conditions: Flow (balanced challenge/skill), easy (insufficient challenge), hard (excessive challenge); EEG recording: Emotiv EPOC X (14 channels, 128 Hz); Session duration: ~30 minutes per condition

\subsubsection{Validation Dataset 2: MultiPENG Video Game Engagement Study}
To assess SIE generalization to interactive gaming contexts, we analyzed the publicly available Multimodal Dataset for Player Engagement Analysis in Video Games (\cite{albuquerque2023multimodal}).

\textbf{Original Study Design }- N = 39 participants; Task: Playing popular video games across varying difficulty levels; EEG recording: [Device specifications to be confirmed from full paper]; Total sessions: 900 annotated gameplay sessions; Psychological metrics: Engagement, interest, stress, and excitement annotations collected during natural game pauses

\subsubsection{Validation Dataset 3: Visual Stimuli EEG Study}
To test whether SIEs occur during passive visual processing, we analyzed the publicly available EEG Dataset for Natural Image Recognition through Visual Stimuli (\cite{faisal2023eeg}).

\textbf{Original Study Design } - N = 32 participants (35 initially recruited; 3 excluded via Vividness of Visual Imagery Questionnaire); Task: Passive viewing of natural object images (apple, flower, car, human face); EEG recording: Emotiv EPOC X (14 channels; AF3, F7, F3, FC5, T7, P7, O1, O2, P8, T8, FC6, F4, F8, AF4); Sampling rate: 128 Hz; Electrode system: International 10-20 system; Hardware filters: Standard Emotiv EPOC X bandpass; Software: EMOTIV Pro for data collection and annotation

\subsubsection{Combined Dataset Summary}
The combined analysis dataset comprised 91 participants across 661 recording sessions, yielding 1,366 detected SIE events. Six independent datasets spanning five cognitive contexts were included: eyes-closed meditation, cognitive flow arithmetic tasks, video game engagement, and passive visual perception. Recordings were collected using three consumer-grade EEG devices—Muse, Emotiv EPOC X, and Emotiv Insight—with the Emotiv EPOC X used across all datasets (Meditation, PhySF, MPENG, VEP), enabling within-device cross-context validation.
% Option 1: Automatically resize to \textwidth with \resizebox
% Requires: \usepackage{graphicx}

\begin{table}[h]
    \centering
    \caption{Dataset Characteristics}
    \label{tab:placeholder_label}
    \resizebox{\textwidth}{!}{%
        \begin{tabular}{lccccc}
            \toprule
            Dataset & Participants & Sessions & SIE Events & Device(s) & Context \\
            \midrule
            Meditation (Files) & 1 & 5 & 38 & EPOC X & Eyes-closed meditation \\
            Meditation (Muse) & 1 & 17 & 165 & Muse & Eyes-closed meditation \\
            Meditation (Insight) & 1 & 12 & 57 & Insight & Eyes-closed meditation \\
            PhySF & 25 & 46 & 447 & EPOC X & Cognitive flow/non-flow tasks \\
            MultiPENG & 36 & 533 & 604 & EPOC X & Video game engagement \\
            VEP & 29 & 48 & 55 & EPOC X & Passive visual perception \\
            \textbf{Combined} & \textbf{91} & \textbf{661} & \textbf{1,366} & \textbf{3 types} & \textbf{5 contexts} \\
            \bottomrule
        \end{tabular}%
    }
\end{table}

\subsubsection{Sample Size Considerations}
This study used a discovery-validation design rather than a hypothesis-testing framework. The discovery phase analyzed all available longitudinal meditation recordings (34 sessions). For validation datasets, participants were included if at least one SIE was detected in their recordings; participants with zero detected events were excluded from analysis.

\textbf{Exclusions by dataset:} PhySF: 1 of 26 original participants excluded (3.8\%); MultiPENG: 3 of 39 excluded (7.7\%); VEP: 3 of 32 excluded (9.4\%). Combined validation sample: 90 of 97 original participants (92.8\%) exhibited at least one SIE.

\textbf{Post-hoc sensitivity analysis:} With 1,366 total events and observed effect sizes (d = 1.44 for peak-based random triplet comparison, p $<$ 0.0001), statistical power exceeded 0.99 for primary comparisons.

\subsection{EEG Data Acquisition and Preprocessing}
Continuous EEG was recorded using three consumer-grade devices: (1) Emotiv EPOC X with 14 channels (AF3, F7, F3, FC5, T7, P7, O1, O2, P8, T8, FC6, F4, F8, AF4) at 128 Hz using saline-soaked felt sensors and CMS/DRL reference; (2) Emotiv Insight with 5 channels (AF3, AF4, T7, T8, Pz) at 128 Hz using semi-dry polymer sensors and CMS/DRL reference; and (3) Muse (2016 model) with 4 channels (TP9, AF7, AF8, TP10) at 256 Hz using dry electrodes with FPz reference. All electrodes were positioned according to the international 10-20 system (\cite{jasper1958tentwenty}). Raw EEG signals were high-pass filtered at 1 Hz using a 4th-order zero-phase Butterworth filter (\cite{butterworth1930filter}) to remove slow drifts while preserving oscillatory activity in the Schumann Resonance frequency range. Data were segmented into continuous recording sessions for subsequent analysis.

\subsection{Schumann Resonance Harmonic Detection}

\subsubsection{Canonical Frequencies}
The Earth's Schumann Resonances arise from electromagnetic standing waves in the Earth-ionosphere cavity, with theoretical frequencies at approximately 7.83 Hz (fundamental) and harmonics near 14.3, 20.8, 27.3, and 33.8 Hz (\cite{schumann1952strahlungslosen}; \cite{nickolaenko2014schumann}). However, measured SR frequencies vary with ionospheric conditions. We derived canonical frequencies and search bandwidths from long-term magnetometer data recorded by the Space Observing System 70 in Tomsk, Russia (https://sos70.ru/?page\_id=48): SR1 = 7.6 Hz (±0.6 Hz), SR2o = 13.75 Hz (±0.75 Hz), SR3 = 20 Hz (±1.0 Hz), SR4 = 25 Hz (±2.0 Hz), and SR5 = 32 Hz (±2.5 Hz). These empirically-derived values served as canonical targets for harmonic detection, with session-specific refinement described below.

\subsubsection{Spectral Parameterization Using FOOOF}
To separate oscillatory (periodic) activity from aperiodic (1/f) background noise, we first estimated power spectra using Welch’s method (\cite{welch1967psd}) and then employed spectral parameterization with the FOOOF algorithm (\cite{haller2018fooof}).
\begin{displaymath}
\log_{10}(PSD(f)) = L(f) + \sum_{i} G_{i}(f)
\end{displaymath}
where L represents the aperiodic component modeled as:
\begin{displaymath}
L(f) = b - \chi \times \log_{10}(f)
\end{displaymath}
% Requires: \usepackage{booktabs}
\begin{table}[h]
    \centering
     \begin{tabular}{llll}
        \toprule
        \multicolumn{4}{c}{\textbf{Table 1: Equation 1 Parameters}} \\
        \midrule
        PSD($f$)& Power spectral density at frequency $f$ & measured & $\mu\mathrm{V}^2/\mathrm{Hz}$ \\
        $L(f)$ & Aperiodic (1/f) component & fitted & $\log_{10}(\mu\mathrm{V}^2/\mathrm{Hz})$ \\
        $G_i(f)$ & $i$-th Gaussian periodic peak & fitted & $\log_{10}(\mu\mathrm{V}^2/\mathrm{Hz})$ \\
        $f$ & Frequency & 1--50 & Hz \\
        \bottomrule
    \end{tabular}
\end{table}
% Requires: \usepackage{booktabs}
\begin{table}[h]
    \centering
    \begin{tabular}{llll}
        \toprule
        \multicolumn{4}{c}{\textbf{Table 2: Equation 2 Parameters}} \\
        \midrule
        $b$ & Aperiodic offset (y-intercept) & fitted & $\log_{10}(\mu\mathrm{V}^2/\mathrm{Hz})$ \\
        $\chi$ & Aperiodic exponent (spectral slope) & 1--2 & dimensionless \\
        $f$ & Frequency & 1--50 & Hz \\
        \bottomrule
    \end{tabular}
\end{table}
% Requires: \usepackage{booktabs}
\begin{table}[h]
    \centering
    \resizebox{\textwidth}{!}{%
    \begin{tabular}{lllll}
        \toprule
        \multicolumn{5}{c}{\textbf{Table 3: FOOOF Fitting Parameters}} \\
        \midrule
        Parameter & Symbol & Value & Units & Rationale \\
        \midrule
        Frequency range & [$f_{\min}, f_{\max}$] & (1.0, 50.0) & Hz & Covers all SR harmonics \\
        Segment length & $T_{\text{seg}}$ & 100.0 & s & Provides 0.01 Hz frequency resolution \\
        Overlap & --- & 0.5 & fraction & 50\% Welch segment overlap \\
        Max peaks & $n_{\max}$ & 7 & count & Sufficient for 5+ SR harmonics \\
        Peak threshold & $\theta_{\text{peak}}$ & 0.1 & $\log_{10}$ units & Lowered for subtle SR peaks \\
        Min peak height & $h_{\min}$ & 0.05 & $\log_{10}$ units & Floor for peak detection \\
        Peak width limits & [$\sigma_{\min}, \sigma_{\max}$] & (0.3, 4.0) & Hz & Accommodates narrow SR peaks \\
        \bottomrule
    \end{tabular}%
    }
\end{table}

\subsubsection{Hybrid Two-Phase Harmonic Detection}
We implemented a hybrid approach combining session-level stability with event-level adaptability:

\begin{itemize}
    \item \textbf{Phase 1 (Session-Level):} Power spectral density was computed across the entire recording session using Welch's method with the parameters specified in Table 3. FOOOF was fitted to extract session-level harmonic frequencies and aperiodic parameters (exponent χ, goodness-of-fit R²).
    \item \textbf{Phase 2 (Per-Event)}: For each detected ignition event (Section 2.3), FOOOF was refitted to the event-specific time window using canonical harmonics as seed frequencies. This captured dynamic frequency shifts during ignition states while maintaining comparability across events.
\end{itemize}

\subsubsection{Peak Matching Strategy}
When multiple spectral peaks occurred within a harmonic search window, we employed maximum power selection to identify the matched frequency. For each canonical harmonic $f_{\text{can}}$, all FOOOF-detected peaks within the search window $[f_{\text{can}} - \Delta f, f_{\text{can}} + \Delta f]$ were identified, and the peak with highest spectral power was selected:
\begin{displaymath}
f_{\text{matched}} = f_k \text{ where } k = \arg\max_i(P_i)
\end{displaymath}
% Requires: \usepackage{booktabs}
\begin{table}[h]
    \centering
    \resizebox{\textwidth}{!}{%
\begin{tabular}{llll}
        \toprule
        Symbol & Definition & Value & Units \\
        \midrule
        $f_i$ & Center frequency of $i$-th candidate peak & detected & Hz \\
        $P_i$ & Power (height) of $i$-th candidate peak & detected & $\log_{10}(\mu\text{V}^2/\text{Hz})$ \\
        \midrule
        \multicolumn{4}{l}{\textit{Canonical frequencies and search half-bandwidths (from Tomsk SOS70):}} \\
        $f_{\text{can,SR1}}$ & SR1 canonical frequency & 7.6 ± 0.6 & Hz \\
        $f_{\text{can,SR2o}}$ & SR2o canonical frequency & 13.75 ± 0.75 & Hz \\
        $f_{\text{can,SR3}}$ & SR3 canonical frequency & 20 ± 1.0 & Hz \\
        $f_{\text{can,SR4}}$ & SR4 canonical frequency & 25 ± 2.0 & Hz \\
        $f_{\text{can,SR5}}$ & SR5 canonical frequency & 32 ± 2.5 & Hz \\
        \bottomrule
    \end{tabular}
%
}
\end{table}

The search window for each canonical frequency was defined as $[f_{\text{can}} - \Delta f_{\text{search}}, f_{\text{can}} + \Delta f_{\text{search}}]$, with harmonic-specific half-bandwidths (Table 4) reflecting greater natural variability at higher SR modes. This maximum power selection approach ensures that the strongest spectral component within each search window is matched to the corresponding SR harmonic, prioritizing signal strength over proximity to canonical frequency.

\subsection{Ignition Event Detection}

We developed a seven-stage detection pipeline to systematically identify, refine, and validate candidate SIE events. The pipeline was designed to balance sensitivity (detecting true SIEs) against specificity (excluding artifacts and non-SIE oscillatory events).

\subsubsection{Schumann Resonance Envelope Computation}
To detect transient increases in Schumann-band activity, we computed a composite SR envelope across all electrodes using the following procedure:

\begin{enumerate}
\item EEG signals from all N = 14 channels were averaged to create a composite signal, reducing channel-specific noise while preserving spatially coherent activity.
\item The composite signal was bandpass filtered at the fundamental frequency (f₀ ± Δf) using a 4th-order Butterworth filter with zero-phase (filtfilt) implementation.
\item The analytic signal was obtained via Hilbert transform (\cite{gabor1946theory}), and the instantaneous envelope was computed as the magnitude of the analytic signal.
\item The envelope was smoothed using a Hanning-windowed moving average of duration T\_smooth.
\item The smoothed envelope was z-score normalized across the entire session:
\end{enumerate}
\begin{displaymath}
z(t) = [env(t) - μ_{env}] / σ_{env}
\end{displaymath}
% Requires: \usepackage{booktabs}
\begin{table}[h]
    \centering
    \begin{tabular}{llll}
        \toprule
        Symbol & Definition & Value & Units \\
        \midrule
        $\mathrm{env}(t)$ & Instantaneous Hilbert envelope magnitude & computed & $\mu$V \\
        $\mu_{\mathrm{env}}$ & Session mean of smoothed envelope & computed & $\mu$V \\
        $\sigma_{\mathrm{env}}$ & Session standard deviation of envelope & computed & $\mu$V \\
        $f_{0}$ & Center frequency (fundamental) & 7.6 & Hz \\
        $\Delta f$ & Half-bandwidth & 0.6 & Hz \\
        $T_{\mathrm{smooth}}$ & Smoothing window duration & 0.01 & s \\
        $N$ & Number of electrodes & 14 & count \\
        $f_{s}$ & Sampling rate & 128 & Hz \\
        \bottomrule
    \end{tabular}
\end{table}

\subsubsection{Onset Detection}
Candidate ignition onsets were identified as rising-edge threshold crossings where $z(t) \geq z_\text{thresh}$.
\begin{table}[h]
    \centering
    \resizebox{\textwidth}{!}{%
    \begin{tabular}{l c c c p{7cm}}
        \hline
        Parameter & Symbol & Value & Units & Description \\
        \hline
        Z-threshold & $z_{\text{thresh}}$ & 3.0 & SD & Detection threshold (p $<$ 0.003 under Gaussian assumptions) \\
        Min ISI & $T_{\text{ISI}}$ & 2.0 & s & Minimum inter-stimulus interval to prevent rapid false positives \\
        Filter order & $n_{\text{order}}$ & 4 & --- & Butterworth filter order \\
        \hline
    \end{tabular}
    %
    }
\end{table}
\subsubsection{Window Definition and Merging}
For each detected onset, an analysis window was defined with the following parameters:
% Requires: \usepackage{booktabs}
\begin{table}[h]
    \centering
    \resizebox{\textwidth}{!}{%
    \begin{tabular}{lllll}
        \toprule
        Parameter       & Symbol    & Value & Units & Description                                    \\
        \midrule
        Window duration & $T_{\text{win}}$   & 20.0  & s     & Total window length                           \\
        Pre-onset       & $T_{\text{pre}}$   & 10.0  & s     & Time before onset (window centered on onset)  \\
        Post-onset      & $T_{\text{post}}$  & 10.0  & s     & Time after onset                              \\
        Merge gap       & $T_{\text{merge}}$ & 10.0  & s     & Maximum gap for merging adjacent windows      \\
        \bottomrule
    \end{tabular}
%
}
\end{table}
Adjacent windows separated by less than T\_merge were merged into single extended events to capture prolonged ignition states. All windows were clipped to recording boundaries.

\subsubsection{Neural Ignition Onset (t₀) Refinement}
The precise timing of neural ignition onset (t₀) was refined using Kuramoto phase synchronization (\cite{kuramoto1984chemical}; \cite{strogatz2000kuramoto}). The Kuramoto order parameter R(t) quantifies instantaneous phase coherence across channels:

\begin{displaymath}
R(t) = \left| \frac{1}{N} \sum_{k} e^{i\phi_k(t)} \right|
\end{displaymath}
The ignition onset t₀ was defined as the time of maximum gradient (dR/dt) where $R(t) \geq R_{\text{thresh}}$, marking the steepest rise in multi-channel phase synchronization—the putative moment of neural ignition.

\textbf{Procedure:}
\begin{enumerate}
\item Bandpass filter all channels at SR1 center frequency ±1.5 Hz (using event-specific SR1 from Stage 2)
\item Compute analytic signal via Hilbert transform
\item Extract instantaneous phase φⱼ(t) for each channel j
\item Compute Kuramoto order parameter:$R(t) = \frac{1}{N}\left|\sum_{j=1}^{N} e^{i\phi_j(t)}\right|$, where N is the number of artifact-free channels.
\item Compute temporal derivative dR/dt using centered finite differences
\item Detect synchronization onset (t₀) at the time of maximum dR/dt where R(t) > 0.6
\item Record R(t₀), maximum R during event, and synchronization onset latency relative to power onset
\end{enumerate}

% Requires: \usepackage{booktabs}
\begin{table}[H]
    \centering
     \resizebox{\textwidth}{!}{%
    \begin{tabular}{llll}
        \toprule
        Symbol & Definition & Value & Units \\
        \midrule
        $R(t)$ & Kuramoto order parameter (phase coherence) & 0--1 & dimensionless \\
        $N$ & Number of EEG channels & 14 & count \\
        $\varphi_k(t)$ & Instantaneous phase of channel $k$ at SR frequency & computed & radians \\
        $i$ & Imaginary unit & $\sqrt{-1}$ & --- \\
        \bottomrule
    \end{tabular}
% 
}
\end{table}

% Requires: \usepackage{booktabs}
\begin{table}[H]
    \centering
    \resizebox{\textwidth}{!}{%
    \begin{tabular}{lllll}
        \toprule
        Parameter       & Symbol        & Value          & Units         & Description                                      \\
        \midrule
        Sliding window  & $T_R$         & 2.5            & s             & $R(t)$ computation window                        \\
        Step size       & $\Delta t_R$  & 0.001          & s             & Temporal resolution (1 ms)                       \\
        R threshold     & $R_{\text{thresh}}$ & 0.6    & dimensionless & Minimum $R$ for $t_0$ detection                  \\
        Edge buffer     & $T_{\text{edge}}$   & 2.0    & s             & Excluded edge samples to avoid boundary artifacts \\
        Min RMS         & $\mathrm{RMS}_{\min}$ & $1\times 10^{-7}$ & $\mu$V & Power threshold for valid $R$ computation        \\
        \bottomrule
    \end{tabular}
%
}
\end{table}

\subsection{Spatio-Spectral Decomposition}

To construct an optimized virtual Schumann Resonance reference signal, we employed Spatio-Spectral Decomposition (\cite{nikulin2011ssd}), a spatial filtering technique that maximizes signal-to-noise ratio at target frequencies.
\begin{displaymath}
C_S · w = λ · C_N · w 
\end{displaymath}

% Requires: \usepackage{booktabs}
\begin{table}[H]
    \centering
    \resizebox{\textwidth}{!}{%
    \begin{tabular}{llll}
        \toprule
        Symbol & Definition & Dimensions & Description \\
        \midrule
        $C_S$ & Signal covariance matrix & $N \times N$ & Covariance of SR-band filtered EEG \\
        $C_N$ & Noise covariance matrix & $N \times N$ & Covariance of flanking-band filtered EEG \\
        $\mathbf{w}$ & Spatial filter weights & $N \times 1$ & Eigenvector solution \\
        $\lambda$ & Eigenvalue & scalar & Signal-to-noise ratio \\
        $N$ & Number of channels & 14 & --- \\
        \bottomrule
    \end{tabular}
%
}
\end{table}
The eigenvector w* corresponding to the maximum eigenvalue λ\_max provides the optimal spatial filter:
\begin{displaymath}
w* = argmax_w (w^T C_S w) / (w^T C_N w) 
\end{displaymath}

% Requires: \usepackage{booktabs}
\begin{table}[H]
    \centering
    \resizebox{\textwidth}{!}{%
    \begin{tabular}{llll}
        \toprule
        Band & Frequency Range & Units & Description \\
        \midrule
        Signal band & [$f_{0}-\Delta f,\, f_{0}+\Delta f$] = [7.33, 8.33] & Hz & SR fundamental $\pm$ 0.5 Hz \\
        Lower flank & [$f_{0}-\Delta f-\Delta f_{\text{flank}},\, f_{0}-\Delta f$] = [6.33, 7.33] & Hz & Below signal band \\
        Upper flank & [$f_{0}+\Delta f,\, f_{0}+\Delta f+\Delta f_{\text{flank}}$] = [8.33, 9.33] & Hz & Above signal band \\
        \bottomrule
    \end{tabular}

%
}
\end{table}
% Requires: \usepackage{booktabs}
\begin{table}[H]
    \centering
    \resizebox{\textwidth}{!}{%
    \begin{tabular}{lllll}
        \toprule
        Parameter & Symbol & Value & Units & Description \\
        \midrule
        Signal bandwidth & $\Delta f$ & 0.4--0.5 & Hz & Half-width of signal band \\
        Flank width & $\Delta f_{\text{flank}}$ & 1.0 & Hz & Width of each flanking noise band \\
        Regularization & $\varepsilon$ & $1\times10^{-12}$ & --- & Pseudoinverse stabilization constant \\
        \bottomrule
    \end{tabular}
%
}
\end{table}
The virtual reference signal was computed as v\_SR(t) = w*\textasciicircum{}T × X(t), where X(t) is the N × T multi-channel EEG matrix. This yields a single-channel reference with enhanced sensitivity to spatially coherent Schumann-band activity while suppressing broadband noise.

\subsection{Per-Event Multi-Scale Characterization}

For each detected ignition event, multiple time-resolved metrics were computed to characterize the dynamics of Schumann-band activity.

\subsubsection{Per-Harmonic SR Envelope Z-Scores}
For each harmonic frequency f\_n, the envelope z-score was computed within the event window using the same procedure as Section 2.3.1.

% Requires: \usepackage{booktabs}
\begin{table}[H]
    \centering
    \resizebox{\textwidth}{!}{%
    \begin{tabular}{llll}
        \toprule
        Metric & Symbol & Window & Description \\
        \midrule
        Peak z-score & $\mathrm{sr}\{n\}\_z\_\mathrm{max}$ & full event & Maximum $z(t)$ in event window \\
        Mean $\pm 5\,\mathrm{s}$ & $\mathrm{sr}\_z\_\mathrm{mean\_pm5}$ & $t_0 \pm 5\,\mathrm{s}$ & Mean z-score around ignition onset \\
        Post-peak mean & $\mathrm{sr}\_z\_\mathrm{mean\_post5}$ & peak $\rightarrow$ $+5\,\mathrm{s}$ & Decay assessment after peak \\
        Peak time & $\mathrm{sr}\_z\_\mathrm{peak\_t}$ & full event & Time of maximum z-score \\
        \bottomrule
    \end{tabular}
%
}
\end{table}

These metrics were computed independently for each harmonic (sr1\_z\_max, sr2\_z\_max, sr3\_z\_max, etc.).

\subsubsection{Per-Event Magnitude Squared Coherence (MSC)}
Mean squared coherence (\cite{carter1973estimation}) was computed for each harmonic frequency using a Hilbert-based approach with median reference:
\begin{displaymath}
MSC = \frac{|S_{xy}(f)|^2}{S_{xx}(f) \cdot S_{yy}(f)}
\end{displaymath}
\textbf{Procedure}:
\begin{enumerate}
\item Compute the cross-spectral density Sxy(f) between each channel and the median reference signal (median voltage across channels at each time point)
\item Compute auto-spectral densities Sxx(f) and Syy(f)
\item Calculate MSC:  \$MSC(f) = \textbackslash{}frac\{|S\_\{xy\}(f)|\textasciicircum{}2\}\{S\_\{xx\}(f) \textbackslash{}cdot S\_\{yy\}(f)\}\$
\item Average MSC across channel pairs within each SR harmonic band (±1 Hz around detected peak)
\item Record mean, maximum, and bandwidth of MSC peak
\end{enumerate}

% Requires: \usepackage{amsmath}
\begin{table}[H]
    \centering
    \resizebox{\textwidth}{!}{%
    \begin{tabular}{llll}
        \hline
        Symbol & Definition & Value & Units \\
        \hline
        $z_{\mathrm{ch}}$ & Analytic signal of channel (Hilbert transform) & computed & complex \\
        $z_{\mathrm{ref}}$ & Analytic signal of median reference across channels & computed & complex \\
        $z_{\mathrm{ref}}^{*}$ & Complex conjugate of reference & computed & complex \\
        $\langle \cdot \rangle$ & Time average over analysis window & --- & --- \\
        \hline
    \end{tabular}
%
}\end{table}

% Requires: \usepackage{booktabs}
\begin{table}[H]
    \centering
    \resizebox{\textwidth}{!}{%
    \begin{tabular}{lllll}
        \toprule
        Parameter & Symbol & Value & Units & Description \\
        \midrule
        Sliding window     & $T_{\mathrm{msc}}$      & 1.0          & s  & MSC computation window \\
        Step size          & $\Delta t_{\mathrm{msc}}$ & 0.1        & s  & 100 ms temporal resolution \\
        Local window       & $T_{\mathrm{local}}$   & $\pm 2.5$    & s  & Window around $t_0$ for peak/mean extraction \\
        Baseline window    & $T_{\mathrm{base\_msc}}$ & $-5$ to $-2$ & s & Pre-event baseline (relative to $t_0$) \\
        Bandpass half-width & $\mathrm{bw}_{\mathrm{msc}}$ & 0.5  & Hz & Per-harmonic filtering \\
        \bottomrule
    \end{tabular}
%
}\end{table}
% Requires: \usepackage{booktabs}
\begin{table}[H]
    \centering
    \begin{tabular}{lll}
        \toprule
        Metric & Symbol & Description \\
        \midrule
        Peak MSC        & \texttt{msc\_peak}     & Maximum MSC in local window \\
        Mean local MSC  & \texttt{msc\_mean\_loc} & Mean MSC in $\pm 2.5\,\mathrm{s}$ around $t_0$ \\
        Baseline MSC    & \texttt{msc\_base}     & Mean MSC in baseline window \\
        MSC AUC         & \texttt{msc\_auc\_loc} & Area under MSC curve in local window \\
        \bottomrule
    \end{tabular}
\end{table}

\subsubsection{Per-Event Phase Locking Value (PLV)}
Phase locking value (\cite{lachaux1999plv}) was computed for each harmonic frequency within a ±5 second window centered on the ignition onset t₀.
\begin{displaymath}
PLV = \frac{1}{T}\left|\sum_{t=1}^{T} e^{i\Delta\phi_{xy}(t)}\right|
\end{displaymath}
\textbf{Procedure:}
\begin{enumerate}
\item Bandpass filter at each SR harmonic frequency (±1.5 Hz)
\item Extract instantaneous phase via Hilbert transform
\item Compute pairwise phase differences Δφxy(t) = φx(t) − φy(t)
\item Calculate PLV: $PLV = \frac{1}{T}\left|\sum_{t=1}^{T} e^{i\Delta\phi_{xy}(t)}\right|$, where T is the number of time points in the event window.\item Average PLV across all channel pairs
\item Record PLV for each SR harmonic
\end{enumerate}
% Requires: \usepackage{array}
\begin{table}[H]
    \centering
    \begin{tabular}{>{\centering\arraybackslash}m{2cm} m{6cm} >{\centering\arraybackslash}m{2cm} >{\centering\arraybackslash}m{2.5cm}}
        \hline
        \textbf{Symbol} & \textbf{Definition} & \textbf{Value} & \textbf{Units} \\
        \hline
        $\Delta\phi(t)$ & Phase difference between channel pairs & computed & radians \\
        $T$ & Number of time samples in analysis window & computed & count \\
        PLV & Phase locking value & 0--1 & dimensionless \\
        $T_{\mathrm{PLV}}$ & Analysis window around $t_{0}$ & $\pm 5$ & s \\
        \hline
    \end{tabular}
\end{table}
\begin{table}[h]
    \centering
    \begin{tabular}{lll}
        \hline
        Metric & Symbol & Description \\
        \hline
        Mean PLV & $sr\{n\}\_{plv}$ & Mean PLV in analysis window \\
        \hline
    \end{tabular}
\end{table}

\subsection{Event Quality Metrics}
Each detected ignition event was characterized by multiple metrics assessing spectral content, cross-frequency relationships, and artifact contamination.

\subsubsection{Harmonic Stack Index (HSI)}
The HSI quantifies the degree to which overtone (harmonic) power dominates over fundamental power. Higher HSI values indicate overtone-dominated events; lower HSI values indicate fundamental-dominated events.

\textbf{Procedure}:
\begin{enumerate}
\item Bandpass filter the signal at the fundamental frequency ($f_0 \pm \Delta f_0$) and compute mean power: $P_f = \langle x_{f_0}^2 \rangle$
\item For each overtone frequency $f_n$ (excluding the fundamental), bandpass filter and compute mean power: $P_n = \langle x_{f_n}^2 \rangle$
\item Sum all overtone powers: $P_{\text{overtones}} = \sum_{n} P_n$
\item Calculate Harmonic Stack Index:
\begin{displaymath}
HSI = \frac{P_{\text{overtones}}}{P_f}
\end{displaymath}
\end{enumerate}

% Requires: \usepackage{array}
\begin{table}[H]
    \centering
    \begin{tabular}{>{\centering\arraybackslash}m{2.5cm} m{6cm} >{\centering\arraybackslash}m{2.5cm} >{\centering\arraybackslash}m{2cm}}
        \hline
        \textbf{Symbol} & \textbf{Definition} & \textbf{Value} & \textbf{Units} \\
        \hline
        $P_f$ & Mean power at fundamental frequency & computed & $\mu \mathrm{V}^2$ \\
        $P_{\text{overtones}}$ & Sum of mean powers at all overtone frequencies & computed & $\mu \mathrm{V}^2$ \\
        $\Delta f_0$ & Fundamental band half-width & 0.5 & Hz \\
        $\Delta f_n$ & Overtone band half-width & 0.5 & Hz \\
        HSI $> 1$ & Overtone-dominated (harmonics stronger than fundamental) & --- & --- \\
        HSI $< 1$ & Fundamental-dominated (fundamental stronger than harmonics) & --- & --- \\
        \hline
    \end{tabular}
\end{table}
\subsubsection{Spectral Slope}
The aperiodic exponent β was computed from the Welch periodogram (3–45 Hz, excluding SR frequencies) via linear regression in log-log space.
% Requires: \usepackage{booktabs}
\begin{table}[H]
    \centering
    \resizebox{\textwidth}{!}{%
    \begin{tabular}{lllll}
        \toprule
        Parameter & Symbol & Value & Units & Interpretation \\
        \midrule
        Frequency range & $[f_{\beta\_\min}, f_{\beta\_\max}]$ & [3, 45] & Hz & Broadband range \\
        Typical neural $\beta$ & $\beta_{\text{neural}}$ & 1--2 & dimensionless & Physiological activity \\
        Artifact indicator & $\beta_{\text{artifact}}$ & $< 0.5$ & dimensionless & Flattened slope suggests broadband contamination \\
        \bottomrule
    \end{tabular}
%
} \end{table}
\subsubsection{Frequency Specificity Index (FSI)}
The FSI measures the concentration of power at SR frequencies relative to total broadband power:

\begin{displaymath}
FSI = P_{SR} / P_{total}
\end{displaymath}

% Requires: \usepackage{booktabs}
\begin{table}[H]
    \centering
    \begin{tabular}{llll}
        \toprule
        Symbol & Definition & Value & Units \\
        \midrule
        $P_{\text{SR}}$ & Summed power within $\pm\Delta f$ of each harmonic & computed & $\mu\text{V}^2$ \\
        $P_{\text{total}}$ & Total power in [3, 45] Hz & computed & $\mu\text{V}^2$ \\
        $\Delta f_{\text{FSI}}$ & Band half-width around each harmonic & 0.5 & Hz \\
        FSI $> 0.35$ & Indicates frequency-specific activity (true ignition) & --- & fraction \\
        FSI $< 0.05$ & Indicates broadband contamination (artifact) & --- & fraction \\
        \bottomrule
    \end{tabular}
\end{table}
\subsection{Event Quality Scoring}
A composite quality score is computed for each ignition event to facilitate ranking and selection, using only the three observed Schumann harmonics (sr1, sr3, sr5):
\begin{displaymath}
\text{SR-Score} = \frac{Z_{\text{weighted}}^{0.7} \times \text{MSC}_{\text{weighted}}^{1.2} \times \text{PLV}_{\text{weighted}}}{1 + \text{HSI}}
\end{displaymath}
% Requires: \usepackage{booktabs}
\begin{table}[H]
    \centering
    \resizebox{\textwidth}{!}{%
    \begin{tabular}{llll}
        \toprule
        Symbol & Definition & Value & Units \\
        \midrule
        $Z_{\text{weighted}}$ & Weighted z-scores across sr1, sr3, sr5 & computed & SD \\
        $\text{MSC}_{\text{weighted}}$ & Weighted coherence across sr1, sr3, sr5 & computed & 0--1 \\
        $\text{PLV}_{\text{weighted}}$ & Weighted PLV across sr1, sr3, sr5 & computed & 0--1 \\
        HSI & Harmonic stack index (penalty for overtone dominance) & computed & ratio \\
        $\alpha$ & Z-score exponent & 0.7 & --- \\
        $\beta$ & MSC exponent & 1.2 & --- \\
        \bottomrule
    \end{tabular}%
} \end{table}
The weighted averaging sums contributions from the three observed Schumann harmonics using fixed weights derived from the golden ratio:
\begin{displaymath}
X_{\text{weighted}} = \sum_{n \in \{1,3,5\}} w_n \times X_{\text{sr}_n}
\end{displaymath}
% Requires: \usepackage{booktabs}
\begin{table}[H]
    \centering
    \caption{Fixed weights for observed Schumann harmonics}
    \begin{tabular}{llll}
        \toprule
        Label & Frequency (Hz) & Weight & Approximate $\varphi^{-e}$ \\
        \midrule
        sr1 & $\sim$7.6 & 0.618 & $\varphi^{-1}$ \\
        sr3 & $\sim$20 & 0.326 & $\varphi^{-2.5}$ \\
        sr5 & $\sim$32 & 0.146 & $\varphi^{-4}$ \\
        \bottomrule
    \end{tabular}
\end{table}
This weighting scheme emphasizes the fundamental frequency (sr1) while progressively attenuating higher harmonic contributions. The HSI penalty in the denominator discourages overtone-dominated events that may reflect non-SR oscillatory activity. By restricting scoring to the three observed Schumann resonance frequencies, this metric enables direct comparison across sessions with different FOOOF-detected harmonic structures.

\subsection{Pipeline Summary}
% Requires: \usepackage{booktabs}
\begin{table}[H]
    \centering
    \resizebox{\textwidth}{!}{%
    \begin{tabular}{llll}
        \toprule
        Stage & Purpose & Key Parameters & Output \\
        \midrule
        1 & Envelope thresholding & $z > 3.0$, duration $\geq$ 3s & Candidate windows \\
        2 & Harmonic refinement & FOOOF, 3 search windows & Peak frequencies \\
        3 & Network synchronization & Kuramoto $R(t)$, $R > 0.6$ & Synchronization onset \\
        4 & Coherence characterization & MSC, PLV per harmonic & Inter-channel coupling \\
        5 & Harmonic organization & HSI computation & Overtone/fundamental ratio \\
        6 & Quality metrics & Slopes, durations & Event characterization \\
        7 & Composite scoring & SR-Score (sr1, sr3, sr5) & Event quality index \\
        \bottomrule
    \end{tabular}%
} \end{table}

\subsection{Advanced SIE Characterization Methods}
Beyond the core detection and quality metrics described above, each Schumann Ignition Event (SIE) was further characterized using advanced signal processing techniques to probe the underlying neural dynamics. These methods assess time-frequency structure, nonlinear phase coupling, cross-frequency interactions, and temporal phase evolution that may distinguish genuine brain-field coupling from spurious oscillatory events. The following subsections detail complementary analyses applied to validated SIE windows.

\subsubsection{Time-Frequency Analysis}
Time-frequency decomposition was performed using Short-Time Fourier Transform (STFT) (\cite{allen1977short}) with the following parameters:
\begin{figure}[H]
    \centering
    \includegraphics[width=1\linewidth]{image_67ce89d9-dc66-486c-af79-f4904aa4b34a.png}
    \caption{Enter Caption}
    \label{fig:placeholder}
\end{figure}
% Requires: \usepackage{booktabs}
\begin{table}[H]
    \centering
    \resizebox{\textwidth}{!}{%
    \begin{tabular}{llll}
        \toprule
        Parameter & Value & Units & Description \\
        \midrule
        Window duration & 2--4 & s & Adaptive based on frequency resolution requirements \\
        Overlap & 75--95 & \% & High overlap for smooth temporal dynamics \\
        Frequency range & 2--40 & Hz & Covers all SR harmonics and gamma \\
        Aggregation & median & --- & Robust to single-channel artifacts \\
        Normalization & row-wise z-score & --- & Highlights temporal dynamics, controls for 1/f decay \\
        \bottomrule
    \end{tabular}%
} \end{table}

\subsubsection{Harmonic Piano Roll Visualization}
Envelope z-scores for each detected SR harmonic were displayed as a heatmap representation:
\begin{figure}[H]
    \centering
    \includegraphics[width=1\linewidth]{image_22dfbaf9-b618-4b89-a675-55a4c1c77989.png}
    \caption{Enter Caption}
    \label{fig:placeholder}
\end{figure}
\begin{itemize}

\item \textbf{Rows}: Individual harmonics (SR1 through SR6 when detectable)
\item \textbf{Columns}: Time samples
\item \textbf{Color scale}: −3 to +3 z-units
\item \textbf{Adaptive display}: Dynamically adjusted to the number of harmonics detected per event
\end{itemize}

\subsubsection{Envelopes and Phase Locking Value (PLV)}
Per-harmonic envelope and PLV time series were computed as described in Sections 2.5.1 and 2.5.3, providing complementary views of amplitude dynamics and inter-channel phase consistency.

\begin{figure}[H]
    \centering
    \includegraphics[width=1\linewidth]{image_ebdaf3e4-ce77-4bed-bf19-d54be4ba8a51.png}
    \caption{Enter Caption}
    \label{fig:placeholder}
\end{figure}

\subsubsection{Harmonic Stack Index (HSI)}
The HSI quantifies the ratio of overtone power to fundamental power, providing a measure of harmonic dominance within each event window. The HSI algorithm comprised the following steps:

\begin{enumerate}
\item \textbf{Fundamental Power}: Bandpass filter at fundamental frequency ($f_0 \pm 0.5$ Hz) and compute mean squared amplitude: $P_f = \langle x_{f_0}^2 \rangle$
\item \textbf{Overtone Powers}: For each overtone frequency $f_n$ (SR2 through SR6, excluding the fundamental), bandpass filter and compute mean power: $P_n = \langle x_{f_n}^2 \rangle$
\item \textbf{Ratio Computation}: $HSI = \frac{\sum_n P_n}{P_f}$
\end{enumerate}
\begin{figure}[H]
    \centering
    \includegraphics[width=1\linewidth]{image_6d63b940-bc31-46e6-86bc-9a064aba9b3a.png}
    \caption{Enter Caption}
    \label{fig:placeholder}
\end{figure}

Higher HSI values indicate overtone-dominated events (harmonics stronger than fundamental); lower HSI values indicate fundamental-dominated events. $\Delta$HSI = HSI(t) $-$ median(HSI) provided baseline-corrected temporal dynamics, smoothed with 450 ms window. In fundamental-led ignition events, we expect HSI to decrease as the fundamental amplitude increases relative to overtones.

\subsubsection{Triadic Bicoherence (Nonlinear Phase Coupling)}
Nonlinear phase coupling among SR harmonic triplets was quantified using bicoherence (\cite{kim1979bispectral}). For frequency triplets $(f_1, f_2, f_3)$ where $f_1 + f_2 \approx f_3$:
\begin{displaymath}
B(f_1, f_2, f_3; t) = \left| \langle \exp[i(\phi_1 + \phi_2 - \phi_3)] \rangle \right|
\end{displaymath}
where φ₁, φ₂, φ₃ are instantaneous phases at the respective frequencies, and ⟨·⟩ denotes averaging over the analysis window.
\begin{figure}[H]
    \centering
    \includegraphics[width=1\linewidth]{image_70fc2741-d2bf-44c7-a12c-88479c1fe793.png}
    \caption{Enter Caption}
    \label{fig:placeholder}
\end{figure}
\textbf{Analyzed Triads:}
\begin{itemize}
\item \textit{Fundamental doubling}: (f₁, f₁, f₂) — tests whether SR1 + SR1 → SR2
\item \textit{Consecutive harmonics}: (f₁, f₂, f₃) — tests sequential coupling
\item \textit{Non-adjacent couplings}: (f₁, f₃, f₄) — tests higher-order relationships
\end{itemize}
% Requires: \usepackage{booktabs}
\begin{table}[H]
    \centering
    
    \begin{tabular}{llll}
        \toprule
        Parameter & Value & Units & Description \\
        \midrule
        Window duration & 0.8 & s & Analysis window \\
        Temporal smoothing & 500 & ms & Smoothing kernel \\
        \bottomrule
    \end{tabular}
    \end{table}

\subsubsection{Theta-Gamma Phase-Amplitude Coupling (PAC)}
Cross-frequency coupling between theta phase and gamma amplitude was quantified using the Mean Vector Length (MVL) method (\cite{canolty2006thetagamma}; \cite{canolty2006gamma}; \cite{tort2010measuring}):
\begin{itemize}
\item \textbf{Phase Signal:} SR fundamental frequency ±1 Hz, averaged across channels
\item \textbf{Amplitude Signal:} Gamma band (25–45 Hz), median across channels
\end{itemize}

\textbf{Procedure}:
\begin{enumerate}
\item Extract instantaneous theta phase via Hilbert transform of bandpass-filtered signal
\item Extract gamma amplitude envelope via Hilbert transform
\item Compute MVL in 6-second sliding windows: $MVL = \frac{1}{N} \left| \sum_{t=1}^{N} A_\gamma(t) \cdot e^{i\phi_\theta(t)} \right|$
\end{enumerate}
\begin{figure}[H]
    \centering
    \includegraphics[width=1\linewidth]{image_d39dfefb-5623-49c1-a96d-91500363c5be.png}
    \caption{Enter Caption}
    \label{fig:placeholder}
\end{figure}
\textbf{Gating}: Analysis restricted to time points in the top 70th percentile of gamma activity to focus on periods of robust high-frequency oscillations.

\subsubsection{Six-Phase Evolution Detection Algorithm}
Six temporal phases were identified using a data-driven segmentation algorithm based on threshold crossings and derivative sign changes:
% Requires: \usepackage{booktabs}
\begin{table}[H]
    \centering
    \resizebox{\textwidth}{!}{%
    \begin{tabular}{lll}
        \toprule
        Phase & Name & Detection Criteria \\
        \midrule
        1 & Baseline & PLV $\leq 1.1\times$ session median AND z-score $< 2.0$ \\
        2 & Coherence-First & PLV rising while power remains near baseline ($t_0$ to $\sim$3s post-onset) \\
        3 & Power Amplification & Rapid power increase (z rising from 2 to peak) \\
        4 & Peak Synchronization & Maximum MSC, PLV, power; $\geq$85\% of peak for $\geq$0.5s \\
        5 & Propagation & PAC surge period ($\sim$5s window around PAC peak) \\
        6 & Decay & Exponential relaxation to baseline \\
        \bottomrule
    \end{tabular}%
} \end{table}

\textbf{Phase Boundary Criteria}:
\begin{itemize}
\item Phase 1→2: PLV exceeds 1.1x baseline median
\item Phase 2→3: Power z-score exceeds 2.0
\item Phase 3→4: Power derivative becomes negative while power remains >85\% of maximum
\item Phase 4→5: Power drops below 85\% of maximum
\item Phase 5→6: PLV drops below 1.5x baseline median
\item Phase 6→1: Both power z-score < 2.0 and PLV ≤ 1.1x baseline median
\end{itemize}

\subsection{Network Dynamics and Connectivity Analysis}
An alternative analytical framework emphasized network connectivity, information flow topology, and global integration dynamics. These methods complement the event-level characterization (Section 2.12) by revealing inter-electrode coupling patterns and cross-scale information transfer during SIEs.

\subsubsection{Harmonic Envelopes}
Per-harmonic amplitude envelopes were computed for each electrode using Hilbert transform of bandpass-filtered signals (±0.5 Hz around each harmonic), providing the input for subsequent connectivity analyses.
\begin{figure}[H]
    \centering
    \includegraphics[width=1\linewidth]{image_c6e6a1c0-bd43-4dae-be40-636653f04f6b.png}
    \caption{Enter Caption}
    \label{fig:placeholder}
\end{figure}
\subsubsection{Frequency-Specific EEG-SR Coherence}
Magnitude-squared coherence between each channel and the SR reference signal (median of all channels) was computed for each harmonic:
\begin{displaymath}
MSC(f) = \frac{|\langle Z_{ch} \cdot \bar{Z}_{ref} \rangle|^2}{\langle |Z_{ch}|^2 \rangle \cdot \langle |Z_{ref}|^2 \rangle}
\end{displaymath}
where Z represents the analytic signal at frequency f.
\begin{figure}[H]
    \centering
    \includegraphics[width=1\linewidth]{image_77afd87f-d10a-4a8b-92ad-46940ca3c0b3.png}
    \caption{Enter Caption}
    \label{fig:placeholder}
\end{figure}

\textbf{Statistical Significance:}
\begin{itemize}
\item \textit{Surrogate testing}: 32 time-shifted reference signals
\item \textit{Threshold}: 95th percentile of surrogate distribution
\end{itemize}

\subsubsection{Kuramoto Global Phase Synchronization Dynamics}
Full timeseries of the Kuramoto order parameter R(t) was computed at the fundamental frequency:
\begin{displaymath}
R(t) = \frac{1}{N} \left| \sum_{j=1}^{N} e^{i\phi_j(t)} \right|
\end{displaymath}
\begin{figure}[H]
    \centering
    \includegraphics[width=1\linewidth]{image_677796f3-6f26-4f90-9b64-6637153a42b4.png}
    \caption{Enter Caption}
    \label{fig:placeholder}
\end{figure}
% Requires: \usepackage{booktabs}
\begin{table}[H]
    \centering
    \begin{tabular}{llll}
        \toprule
        Parameter & Value & Units & Description \\
        \midrule
        Temporal smoothing & 300 & ms & Gaussian kernel \\
        Phase overlay & 6 phases & --- & Color-coded temporal spans \\
        \bottomrule
    \end{tabular}
\end{table}
\subsubsection{Power-Coherence Phase Space Analysis}
Two-dimensional phase portraits were constructed:
\begin{itemize}
\item \textit{X-axis}: SR envelope power (log scale)
\item \textit{Y-axis}: Kuramoto R
\item \textit{Color coding}: Blue = baseline epochs, Orange = ignition epochs
\end{itemize}
\begin{figure}[H]
    \centering
    \includegraphics[width=1\linewidth]{image_a37b2663-b4e5-494e-a5b3-e079c69e051d.png}
    \caption{Enter Caption}
    \label{fig:placeholder}
\end{figure}
This visualization reveals:
\begin{itemize}
\item Power-law relationships between power and synchronization
\item Threshold dynamics (critical R values for state transitions)
\item Hysteresis effects (different trajectories for ignition vs. decay)
\end{itemize}

\subsubsection{Directed Information Flow Topology (Transfer Entropy Proxy)}
Directed information flow between electrode pairs was quantified using a transfer entropy (\cite{schreiber2000transfer}) proxy based on lagged correlations:
\begin{displaymath}
TE(i \rightarrow j) = \rho^2(X_i[t], X_j[t+\tau]) - \rho^2(X_j[t], X_j[t+\tau])
\end{displaymath}
where τ = 50 ms (lag), and ρ² is squared Pearson correlation.
\begin{figure}[H]
    \centering
    \includegraphics[width=1\linewidth]{image_2880ad17-df5f-415d-8e1f-969a9ce7ae07.png}
    \caption{Enter Caption}
    \label{fig:placeholder}
\end{figure}
\textbf{Baseline Correction: }$\Delta TE = TE_{ignition} - TE_{baseline}$. Both terms z-scored within session. Positive ΔTE indicates increased directed coupling during SIEs.
% Requires: \usepackage{booktabs}
\begin{table}[H]
    \centering
    \begin{tabular}{llll}
        \toprule
        Parameter & Value & Units & Description \\
        \midrule
        Lag ($\tau$) & 50 & ms & Temporal offset \\
        Correlation & $\rho^2$ & --- & Squared Pearson \\
        \bottomrule
    \end{tabular}
\end{table}
\subsubsection{Cross-Scale Information Transfer (Theta→Gamma)}
Transfer entropy between theta and gamma envelopes was computed to quantify cross-frequency information flow:
\begin{itemize}
    \item \textbf{Theta envelope:} Fundamental frequency ±1 Hz, median across channels
    \item \textbf{Gamma envelope:} 30–60 Hz, median across channels
\end{itemize}
\begin{displaymath}
TE_{\theta \rightarrow \gamma} = \rho^2(\theta_{env}[t], \gamma_{env}[t+100ms]) - \rho^2(\gamma_{env}[t], \gamma_{env}[t+100ms])
\end{displaymath}
\begin{figure}[H]
    \centering
    \includegraphics[width=1\linewidth]{image_bb07aa05-e696-4e81-b181-098d55d94f29.png}
    \caption{Enter Caption}
    \label{fig:placeholder}
\end{figure}
% Requires: \usepackage{booktabs}
\begin{table}[H]
    \centering
    \begin{tabular}{llll}
        \toprule
        Parameter & Value & Units & Description \\
        \midrule
        Lag & 100 & ms & Theta-to-gamma delay \\
        Window duration & 1.0 & s & Sliding windows \\
        Temporal smoothing & 300 & ms & Smoothing kernel \\
        \bottomrule
    \end{tabular}
\end{table}

\subsection{Null Control Analyses}
Six null control analyses were implemented to validate that SIE characteristics represent genuine neural organization rather than methodological artifacts, random spectral coincidences, or statistical flukes.

\subsubsection{Null Control A: Phase Randomization}
Test whether high synchronization values (MSC, PLV, SR-Score) reflect genuine inter-regional phase coupling or arise from spectral content alone.

\textbf{Procedure}:
\begin{enumerate}
\item For each detected SIE, extract the multi-channel EEG segment
\item Compute the Fast Fourier Transform (FFT) of each channel
\item Randomize phases by replacing each phase value with a random value drawn uniformly from [0, 2π], while preserving the amplitude spectrum exactly
\item Apply inverse FFT to generate phase-shuffled surrogate
\item Repeat steps 2–4 to generate 100 surrogates per event
\item Recompute SR-Score, MSC, and PLV for each surrogate
\item Compare observed values to surrogate distributions
\end{enumerate}

\textbf{Statistical Assessment:}
\begin{itemize}
\item Z-score of observed value relative to surrogate distribution
\item Percentile rank of observed value
\item Cohen's d comparing observed versus surrogate means
\end{itemize}

\subsubsection{Null Control B: Random Temporal Windows (Baseline Comparison)}
Test whether SR-frequency organization is specific to detected SIE events or exists continuously throughout recordings.

\textbf{Procedure}:
\begin{enumerate}
\item Identify baseline periods: segments >10 seconds from any detected SIE
\item Randomly sample 3-second windows from baseline periods (n = 1,366, matching SIE count)
\item Apply identical FOOOF-based peak extraction within SR harmonic search windows
\item Compare baseline windows to SIE windows on:   - Peak detection rate (proportion with detectable peaks in each band)   - Peak frequencies (if detected)   - Frequency ratios   - Power z-scores   - MSC and PLV values
\end{enumerate}

\textbf{Interpretation}: If SR frequencies exist only during SIEs, baseline windows should show no SR-band peaks. If SR frequencies are continuous, baseline windows should show peaks but with different quality metrics.

\subsubsection{Null Control C: Random Spectral Peak Triplets (Uniform)}
Test whether golden ratio precision of SR harmonic frequency ratios is unique or could arise from randomly selected FOOOF peak triplets. This uses uniform sampling from physiologically plausible frequency ranges.

\textbf{Procedure}:
\begin{enumerate}
\item Generate 10,000 random frequency triplets by sampling three frequencies uniformly from the 5–40 Hz range
\item For each triplet (f₁, f₂, f₃), compute the three frequency ratios: f₂/f₁, f₃/f₁, f₃/f₂
\item Compute the deviation from golden ratio predictions:   - Predicted ratios: φ (1.618), φ² (2.618), φ³ (4.236)   - Deviation = mean absolute difference between observed and closest φⁿ value
\item Compare the distribution of random triplet deviations to observed SIE deviation
\end{enumerate}

\textbf{Statistical Assessment:}
\begin{itemize}
\item Percentile rank of observed deviation within random distribution
\item Cohen's d between observed deviation and random mean
\item Probability of obtaining observed or better precision by chance
\end{itemize}

\textbf{Limitation}: Uniform sampling from ranges centered on $\phi^n$ frequencies (e.g., 5--10 Hz centered on $\phi^0 \approx 7.6$ Hz) may bias the null toward producing decent $\phi$-ratios. Null Control~D addresses this limitation.

\subsubsection{Null Control D: Peak-Based Random Triplets}
Test whether SIE frequency triplets show better φ-convergence than random triplets sampled from the actual distribution of FOOOF peaks in EEG data. This controls for the spectral structure of real EEG recordings.

\textbf{Procedure}:
\begin{enumerate}
\item Pool all FOOOF-detected peaks from the complete dataset (N = 244,955 peaks from 661 sessions, spanning all cognitive contexts and devices)
\item Filter peaks into SR harmonic bands: SR1 (6.5–9.5 Hz, n = 23,075), SR3 (18–22 Hz, n = 27,659), SR5 (30–36 Hz, n = 18,347)
\item Generate 10,000 random triplets by independently sampling one peak frequency from each band pool
\item Compute composite φ-error for each triplet using the same metric as SIE events
\item Compare SIE φ-error distribution to random peak triplet distribution
\end{enumerate}

\textbf{Statistical Assessment:}
\begin{itemize}
\item Permutation test (10,000 iterations)
\item Mann-Whitney U test (non-parametric)
\item Cohen's d effect size
\item Bootstrap 95\% CI for SIE mean φ-error
\end{itemize}

\textbf{Interpretation}: Unlike the uniform null (Null Control C), the peak-based null directly tests: ``Given that EEG spectra naturally contain peaks at various frequencies, do SIE events select frequency combinations that are more φ-aligned than typical?'' A significant result indicates that SIE detection preferentially identifies genuinely exceptional φⁿ organization.

\subsubsection{Null Control E: Shuffled Bootstrap (Within-Event Coordination)}
Test whether harmonic frequency relationships arise from event-level coordination (frequencies within an event are coupled) or from population-level marginal distributions (each frequency follows an independent distribution that happens to produce good ratios).

\textbf{Procedure}:
\begin{enumerate}
\item Extract the three harmonic frequencies (SR1, SR3, SR5) from all 1,366 events
\item Create shuffled datasets by independently permuting each harmonic frequency vector:   - Shuffle SR1 frequencies across events (random assignment)   - Shuffle SR3 frequencies across events (independent of SR1 shuffle)   - Shuffle SR5 frequencies across events (independent of SR1, SR3 shuffles)
\item This preserves the marginal distribution of each frequency while destroying within-event pairings
\item Compute the composite φ-error (mean absolute deviation from golden ratio predictions across all three ratios) for each shuffled dataset
\item Repeat 1,000 times to generate null distribution
\item Compare observed composite φ-error to null distribution
\end{enumerate}

\textbf{Interpretation}: If observed error is significantly lower than shuffled error, within-event frequency coordination contributes to precision beyond marginal distributions alone.

\subsubsection{Null Control F: Distributional Null Model}
Quantify the expected precision arising from the marginal frequency distributions alone, providing a more precise baseline for evaluating within-event effects.

\textbf{Procedure}:
\begin{enumerate}
\item Model each harmonic frequency as a univariate normal distribution using observed means and standard deviations
\item Generate 10,000 synthetic datasets by independently sampling from each distribution
\item Compute frequency ratios and φ-error for each synthetic dataset
\item Compare observed error to the distribution of synthetic errors
\end{enumerate}

\textbf{Interpretation}: Similarity between observed and synthetic precision suggests that marginal distributions largely determine ratio precision; substantial improvement suggests additional within-event coordination.

\subsection{Cross-Validation Analyses}

\subsubsection{Device Independence}
To test whether SIE characteristics are device-independent, one-way ANOVA was performed comparing events from Muse (n = 165), Emotiv EPOC X (n = 1,144), and Emotiv Insight (n = 57) devices on all primary metrics (harmonic frequencies, MSC, PLV, SR-Score).

\subsubsection{Context Independence}
To test whether SIEs differ between meditation and cognitive task contexts, independent-samples t-tests compared meditation (n = 260) versus PhySF (n = 447) events on all primary metrics.

\subsubsection{Task Condition Comparison}
Within the PhySF dataset, paired-samples t-tests compared flow versus non-flow conditions on SIE detection rate (events per session) and event characteristics.

\subsection{Golden Ratio Analysis}
An unexpected observation during initial characterization was that SIE harmonic frequencies approximate powers of the golden ratio (φ = 1.6180339...) (\cite{livio2002golden}). To quantify this relationship:

\subsubsection{Predicted Frequencies}
Analysis of detected SIE events reveals that harmonic frequencies follow powers of the golden ratio scaled by the fundamental frequency:
\begin{equation}
    f_n = f_0 \cdot \varphi^n
\end{equation}
where $f_0 = 7.6$ Hz is the Schumann fundamental frequency determined from SR observatory data, and $n$ is an integer representing harmonic level.

This empirical relationship can be connected to physical theory by assuming $f_0 = \frac{c}{2\pi r}$, where $c$ is the speed of light (299,792,458 m/s) and $r$ is Earth's radius (6,367,000 m). Converting to angular frequency $\omega = 2\pi f$:
\begin{equation}
    \omega_n = 2\pi f_0 \cdot \varphi^n = 2\pi \cdot \frac{c}{2\pi r} \cdot \varphi^n = \frac{c}{r} \cdot \varphi^n
\end{equation}

Using $f_0 = 7.6$ Hz, the predicted neural frequencies are:
% Requires: \usepackage{booktabs}
\begin{table}[H]
    \centering

    \begin{tabular}{lll}
        \toprule
        $n$ & $\varphi^n$ & Predicted Frequency (Hz) \\
        \midrule
        0 & 1.000 & 7.60 (fundamental) \\
        2 & 2.618 & 19.90 (SR3) \\
        3 & 4.236 & 32.19 (SR5) \\
        \bottomrule
    \end{tabular}
\end{table}
\subsubsection{Ratio Precision Analysis}

For each SIE event, three frequency ratios were computed:
\begin{itemize}
\item SR3/SR1 (predicted: φ² = 2.618)
\item SR5/SR1 (predicted: φ³ = 4.236)
\item SR5/SR3 (predicted: φ = 1.618)
\end{itemize}

Percentage error was computed as: |observed − predicted| / predicted × 100

\subsubsection{Independence-Convergence Analysis}
To test whether individual frequencies vary independently while ratios converge, we computed:
\begin{enumerate}
\item Pairwise Pearson correlations between SR1, SR3, and SR5 frequencies across events
\item The coefficient of variation (CV) for each frequency and each ratio
\item Scatter plots of frequency pairs with theoretical prediction lines
\end{enumerate}

\subsection{Statistical Analysis}

\subsubsection{General Approach}
All statistical tests were two-tailed with α = 0.05. Effect sizes are reported as Cohen's d for t-tests (small: 0.2, medium: 0.5, large: 0.8) and Pearson's r for correlations.

\subsubsection{Multiple Comparisons}
For analyses involving multiple comparisons (e.g., three harmonic frequencies × multiple metrics), Bonferroni correction was applied where indicated.

\subsubsection{Non-Parametric Alternatives}
When distributional assumptions were violated (Shapiro-Wilk p < 0.05), non-parametric alternatives were employed: Mann-Whitney U for independent samples, Wilcoxon signed-rank for paired samples, Kruskal-Wallis for multiple groups.

\subsubsection{Bootstrap Confidence Intervals}
For derived metrics without known sampling distributions (e.g., composite φ-error), 95\% confidence intervals were estimated via percentile bootstrap with 10,000 resamples.

\subsection{Data and Code Availability}

\subsubsection{Meditation Dataset}
Longitudinal meditation recordings will be made available upon publication via:
\begin{itemize}
    \item https://github.com/neurokinetikz/schumann.
\end{itemize}

\subsubsection{Analysis Code}
Complete Python analysis code, including the SIE detection pipeline, null control implementations, and statistical analyses, will be made available upon publication at:
\begin{itemize}
    \item https://github.com/neurokinetikz/schumann.
\end{itemize}

\subsubsection{Derived Data}
Event-level data for all 1,366 detected SIEs, including harmonic frequencies, synchronization metrics, and quality scores, will be provided as supplementary material.

\subsection{Declarations}

\subsubsection{Conflicts of Interest}
The authors declare no competing financial interests or personal relationships that could have appeared to influence the work reported in this paper.

\subsubsection{Funding}
This research did not receive any specific grant from funding agencies in the public, commercial, or not-for-profit sectors.

\subsection{Software and Computing Environment}
All analyses were performed on a MacBook Air (Apple M4, 16 GB RAM) running macOS Tahoe 26.1 using the following software environment:

% Requires: \usepackage{booktabs}
\begin{table}[H]
    \centering
    \resizebox{\textwidth}{!}{%
    \begin{tabular}{lll}
        \toprule
        Software & Version & Purpose \\
        \midrule
        Python & 3.9.7 & Primary analysis environment \\
        NumPy (\cite{harris2020numpy}) & 1.21.0 & Numerical computing \\
        SciPy (\cite{virtanen2020scipy}) & 1.7.1 & Scientific computing, signal processing \\
        MNE-Python (\cite{gramfort2013mne}) & 0.24.0 & EEG data handling, spectral analysis \\
        FOOOF (\cite{donoghue2020fooof}) & 1.0.0 & Spectral parameterization \\
        Matplotlib (\cite{hunter2007matplotlib})  & 3.4.3 & Visualization \\
        Pandas (\cite{mckinney2010pandas}) & 1.3.3 & Data management \\
        \bottomrule
    \end{tabular} %
} \end{table}

\section{Results}
\subsection{Dataset Overview and Event Detection}
Application of the seven-stage SIE detection pipeline to the combined dataset yielded 1,366 discrete events across 661 recording sessions from 91 unique participants (Table 2). Events were detected across all six datasets and five cognitive contexts, with overall mean detection rate of 2.1 ± 3.3 events per session (range: 1–27; individual dataset means ranged from 1.1 to 9.7 events per session). Event duration averaged 26.9 ± 21.3 seconds (median: 20.0 seconds; minimum window size: 20 seconds), consistent with the transient, high-coherence nature of SIE phenomena and characteristic timescales of transient oscillatory events (\cite{siegel2012spectral}).
% Requires: \usepackage{booktabs}
\begin{table}[H]
    \centering
    \resizebox{\textwidth}{!}{%
    \begin{tabular}{lllllll}
        \toprule
        Dataset & Participants & Sessions & SIE Events & Events/Session & Device(s) & Context \\
        \midrule
        PhySF & 25 & 46 & 447 & $9.7 \pm 6.2$ & EPOC X & Flow/Non-flow \\
        Meditation & 1 & 34 & 260 & $7.6 \pm 4.8$ & Muse, EPOC X, Insight & Eyes-closed \\
        MultiPENG & 36 & 533 & 604 & $1.1 \pm 0.4$ & EPOC X & Gaming \\
        VEP & 29 & 48 & 55 & $1.1 \pm 0.5$ & EPOC X & Visual \\
        \midrule
        \textbf{Combined} & \textbf{91} & \textbf{661} & \textbf{1,366} & $\mathbf{2.1 \pm 3.3}$ & \textbf{3 types} & \textbf{5 contexts} \\
        \bottomrule
    \end{tabular}%
} \end{table}
The distribution of events across datasets revealed that PhySF cognitive flow recordings and meditation recordings yielded the highest per-session detection rates (9.7 and 7.6 events/session respectively), with gaming and visual perception tasks showing lower rates (\textasciitilde{}1 event/session). This pattern may reflect differential engagement of SR-aligned neural mechanisms across cognitive states, though recording duration differences may also contribute.

\subsection{Harmonic Definitions}

Spectral analysis of detected SIEs revealed nine distinct harmonic frequencies spanning the 7–41 Hz range (Table 4). These harmonics, labeled SR1 through SR6 (with half-integer labels SR1.5, SR2.5, and quarter-integer SR2o), serve as the basis for all subsequent analyses. The fundamental frequency (SR1) centers near 7.6 Hz at the theta/alpha boundary, with higher harmonics extending through the alpha, beta, and gamma bands.

\textbf{Table 4. Harmonic Frequency Definitions}
% Requires: \usepackage{booktabs}
\begin{table}[H]
    \centering
    \begin{tabular}{lllll}
        \toprule
        Harmonic & Measured (Hz) & CV (\%) & $N$ & EEG Band \\
        \midrule
        SR1 & $7.63 \pm 0.33$ & 3.18 & 1121 & $\theta/\alpha$ Boundary \\
        SR1.5 & $9.96 \pm 0.35$ & 3.56 & 1139 & $\alpha$ Attractor \\
        SR2 & $11.98 \pm 0.40$ & 3.37 & 1050 & $\alpha/\beta$ Boundary \\
        SR2o & $13.76 \pm 0.44$ & 3.17 & 1150 & Low $\beta$ \\
        SR2.5 & $15.50 \pm 0.46$ & 2.99 & 1188 & $\beta$ Attractor \\
        SR3 & $19.99 \pm 0.57$ & 2.02 & 1250 & $\beta$ Boundary \\
        SR4 & $24.98 \pm 1.17$ & 4.69 & 1359 & High $\beta$ Attractor \\
        SR5 & $32.57 \pm 1.23$ & 1.79 & 1364 & $\beta/\gamma$ Boundary \\
        SR6 & $40.11 \pm 1.61$ & 4.01 & 1365 & $\gamma$ Attractor \\
        \bottomrule
    \end{tabular}
\end{table}
CV = coefficient of variation. N = events with detectable harmonic. EEG band assignments follow canonical frequency definitions (\cite{buzsaki2004neuronal}).

\subsection{Exemplar Event Walkthrough}
To illustrate these harmonics in action, we present a detailed analysis of a canonical high-quality SIE event before examining population-level statistics.

Figure 3 presents a canonical high-quality SIE event (SR-score = 1.04, top 1.5\% of events) recorded during eyes-closed meditation using the Emotiv EPOC X device. This 26-second event (556–582s) illustrates the complete temporal dynamics and multi-scale organization characteristic of SIE phenomena. Each panel reveals distinct aspects of the ignition process. It should be noted that while this exemplar exhibits the general characteristics of SIE phenomena, individual events vary considerably in their specific dynamics, timing, and magnitude of effects; the population-level statistics presented subsequently capture this variability.

\subsubsection{Spectral and Temporal Dynamics}
\begin{figure}[H]
    \centering
    \includegraphics[width=1\linewidth]{download (5).png}
    \caption{Exemplar SIE event spectral and temporal dynamics. (A) Time-frequency spectrogram showing harmonic activation at φⁿ frequencies. (B) Harmonic piano-roll showing z-score elevation across all six harmonics. (C) Envelope traces overlaid with phase-locking value, demonstrating coherence-first dynamics. (D) Harmonic Stack Index showing fundamental dominance (decreasing HSI) preceding ignition.}
    \label{fig:exemplar_spectral}
\end{figure}

\textbf{Panel A: SR Spectrogram (2–40 Hz).}
The time-frequency representation reveals broadband spectral changes during ignition. Horizontal dashed lines mark the predicted φⁿ harmonic frequencies (7.72, 12.03, 13.86, 20.28, 24.92, 32.44 Hz). Key observations:
\begin{itemize}
\item \textbf{Pre-ignition (556–564s):} Diffuse spectral activity without clear harmonic organization. Power distributed across frequencies with no preferential concentration at φⁿ positions.
\item \textbf{Coherence/Plateau phase (564–568s):} Emergence of distinct horizontal "bands" of elevated power precisely aligned to the dashed harmonic markers. This alignment confirms that power increases occur specifically at φⁿ frequencies rather than broadband.
\item \textbf{Ignition (568–570s):} Maximal power concentration at harmonic frequencies, visible as bright yellow horizontal striations. The fundamental (7.72 Hz) and second harmonic (12.03 Hz) show strongest activation.
\item \textbf{Propagation (570–578s):} Sustained but gradually diminishing harmonic structure. Higher harmonics (20–32 Hz) maintain elevation longer than would be expected from simple theta-band filtering artifacts.
\end{itemize}

\textbf{Panel B: Harmonic Piano-Roll (envelope z-score).}
The piano-roll representation isolates activity at each φⁿ harmonic, with z-score color mapping (−3 to +3σ) referenced to session baseline. Y-axis labels show harmonic order and measured frequency: 1× (7.72 Hz), 2× (12.03 Hz), 3× (13.86 Hz), 4× (20.28 Hz), 5× (24.92 Hz), 6× (32.44 Hz).

Critical observations:
\begin{itemize}
\item \textbf{Multi-band activation:} All six harmonics show simultaneous z > 2σ elevation during ignition, confirming this is not merely theta-band activity bleeding into higher frequencies.
\item \textbf{Activation sequence:} The fundamental (1×) activates first (\textasciitilde{}565s), followed by sequential activation of higher harmonics over \textasciitilde{}2 seconds. This bottom-up cascade suggests hierarchical harmonic generation rather than independent oscillators.
\item \textbf{Differential persistence:} Lower harmonics (1×, 2×) show sustained elevation throughout propagation (570–578s), while higher harmonics (5×, 6×) return to baseline more rapidly. This temporal dissociation argues against simple spectral leakage.
\item \textbf{Peak z-scores:} Fundamental shows z > 3 (yellow), while $6\times$ harmonic shows $z \approx 1.5$--$2$ (green-yellow), consistent with the population gradient (SR1: 97.8\% > 3σ; SR6: 40.7\% > 3σ).
\end{itemize}

\textbf{Panel C: Envelopes and PLV.}
This panel overlays z-score envelopes at each harmonic (solid colored lines) with phase-locking value at the fundamental frequency (blue dashed line, right y-axis). The six-phase temporal structure is annotated at bottom.

\textit{The coherence-first signature:}
\begin{itemize}
\item \textbf{PLV rises first:} Phase-locking begins increasing at \textasciitilde{}562s, reaching peak (\textasciitilde{}0.82) at \textasciitilde{}566s—approximately 2–3 seconds BEFORE z-score peaks.
\item \textbf{Z-scores follow:} Power envelopes remain near baseline until \textasciitilde{}565s, then rise sharply to peak at \textasciitilde{}568s.
\item \textbf{Temporal offset:} This \textasciitilde{}2–3 second lag between coherence and amplitude is a hallmark of network-level coordination, inconsistent with local oscillatory bursts where amplitude and phase would co-vary.
\end{itemize}

\textbf{Panel D: Fundamental Dominance (ΔHSI).}
The Harmonic Stack Index deviation from session median tracks the ratio of overtone power to fundamental power. Negative ΔHSI indicates the fundamental is stronger relative to overtones than typical (fundamental-dominated); positive ΔHSI indicates overtone-dominated activity.

Key pattern:
\begin{itemize}
\item \textbf{Baseline:} ΔHSI fluctuates around 0 to +0.05 (typical overtone/fundamental balance).
\item \textbf{Pre-ignition shift:} ΔHSI begins decreasing at \textasciitilde{}563s, reaching minimum (−0.10) during the plateau phase—BEFORE peak power elevation. This indicates the fundamental (SR1) is gaining power relative to overtones.
\item \textbf{Fundamental dominance precedes ignition:} The shift toward fundamental-led activity is established \textasciitilde{}3 seconds before maximal amplitude, suggesting that SR1 dominance is a prerequisite for, not consequence of, ignition.
\item \textbf{Relaxation during propagation:} ΔHSI returns toward zero during propagation, indicating gradual return to typical overtone/fundamental balance.
\end{itemize}

This temporal sequence (fundamental dominance → coherence → amplitude) strongly supports a model where SR1-led frequency organization must be established before power can increase.

\subsubsection{Cross-Frequency Coupling}

\textbf{Panel E: Bicoherence (SR Triads).}
Bicoherence quantifies nonlinear (quadratic) phase coupling between frequency triplets. Each trace represents a specific triad: (f₁, f₂ → f₁+f₂). Legends show triads including (7.72, 7.72→12.03), (7.72, 12.03→13.86), (7.72, 12.03→20.28), etc.

Observations:
\begin{itemize}
\item \textbf{Baseline bicoherence:} \textasciitilde{}0.2–0.4 across triads, indicating modest nonlinear coupling during normal activity.
\item \textbf{Ignition-phase surge:} Multiple triads simultaneously increase to bicoherence > 0.8 during the coherence/plateau phases (\textasciitilde{}565–570s). This coordination across triads indicates that nonlinear harmonic generation is not frequency-specific but affects the entire φⁿ stack.
\item \textbf{Triad synchrony:} The remarkable synchronization of bicoherence increases across different frequency combinations suggests a common generative mechanism rather than independent oscillator coupling.
\item \textbf{Peak timing:} Bicoherence peaks slightly BEFORE maximal z-scores, consistent with nonlinear coupling driving harmonic power rather than vice versa.
\end{itemize}

\textbf{Panel F: θ→γ Phase-Amplitude Coupling (PAC).}
Mean Vector Length (MVL) quantifies the modulation of gamma-band (30–45 Hz) amplitude by theta-band (4–8 Hz) phase—a hallmark of hierarchical neural coding (\cite{lisman2013thetagamma}). Higher MVL indicates stronger cross-frequency coupling.

Temporal profile:
\begin{itemize}
\item \textbf{Baseline (556–563s):} MVL \textasciitilde{}0.05–0.08, indicating weak theta-gamma coupling.
\item \textbf{Pre-ignition rise (563–566s):} MVL begins increasing, reaching \textasciitilde{}0.15 before z-score elevation.
\item \textbf{Peak PAC (566–568s):} MVL reaches maximum (\textasciitilde{}0.20) during the plateau phase, representing a \textasciitilde{}3× increase from baseline.
\item \textbf{Sustained coupling (568–577s):} MVL remains elevated (\textasciitilde{}0.12–0.15) throughout propagation.
\item \textbf{Decay (577–582s):} Gradual return to baseline levels.
\end{itemize}

The theta-gamma coupling increase precedes and accompanies ignition, suggesting that cross-frequency phase relationships are integral to SIE dynamics. The \textasciitilde{}3× increase in MVL during ignition indicates substantial enhancement of hierarchical frequency organization.

\subsubsection{Network Synchronization}
\begin{figure}[H]
    \centering
    \includegraphics[width=1\linewidth]{download (6).png}
    \caption{Exemplar SIE event network dynamics. (G) Harmonic envelope comparison between baseline and ignition states. (H) EEG-SR coherence by electrode and harmonic. (I) State-space trajectory showing coherence vs. power. (J) Kuramoto order parameter time series. (K) Transfer entropy matrix. (L) Cross-scale theta-gamma correlation dynamics.}
    \label{fig:exemplar_network}
\end{figure}

\textbf{Panel G: Harmonic Envelopes (Baseline vs. Ignition).}
Bar chart comparing mean envelope magnitude during baseline (blue) versus ignition (orange) periods at each harmonic frequency.
\begin{table}[H]
    \centering
    \begin{tabular}{llll}
        \toprule
        Harmonic & Baseline & Ignition & Ratio \\
        \midrule
        7.72 Hz (SR1) & $\sim$2.4 & $\sim$4.5 & \textbf{1.88$\times$} \\
        12.03 Hz (SR2) & $\sim$1.4 & $\sim$1.5 & 1.07$\times$ \\
        13.86 Hz (SR2o) & $\sim$1.2 & $\sim$1.5 & 1.25$\times$ \\
        20.28 Hz (SR3) & $\sim$1.2 & $\sim$1.3 & 1.08$\times$ \\
        24.92 Hz (SR4) & $\sim$0.7 & $\sim$0.8 & 1.14$\times$ \\
        32.44 Hz (SR5) & $\sim$0.6 & $\sim$0.7 & 1.17$\times$ \\
        \bottomrule
    \end{tabular} \end{table}
The fundamental frequency shows dramatically greater ignition/baseline ratio (1.88×) compared to higher harmonics (1.07–1.25×), confirming that SR1 is the primary driver of ignition events.

\textbf{Panel H: EEG-SR Coherence at Harmonics (by Electrode).}
Bar chart showing magnitude-squared coherence (MSC) between each of 14 EEG channels and a synthetic SR reference signal, computed separately for each harmonic frequency. All 14 electrodes show MSC > 0.8 at SR1 (7.72 Hz), indicating whole-brain synchronization at the fundamental frequency. MSC systematically decreases at higher harmonics, from \textasciitilde{}0.85–0.95 at SR1 to \textasciitilde{}0.30–0.50 at SR5. All electrode regions (frontal, temporal, parietal, occipital) show similar patterns, arguing against localized generators.

\textbf{Panel I: Coherence vs. SR Power State Space.}
Scatter plot showing the trajectory through coherence (Kuramoto R, y-axis) versus SR envelope power (x-axis, log scale) during baseline (blue) and ignition (orange). Baseline points cluster at low power (\textasciitilde{}1–5) and moderate coherence (R \textasciitilde{}0.5–0.7). The ignition trajectory extends rightward to power >100 while maintaining or increasing coherence (R \textasciitilde{}0.7–0.95), forming a characteristic "lasso" pattern. The clear non-overlap between baseline and ignition states demonstrates that SIEs represent a distinct neural regime.

\textbf{Panel J: Kuramoto Global Integration R(t).}
Time series of the Kuramoto order parameter R(t), measuring instantaneous global phase synchronization across all electrodes. The six-phase temporal structure is visible:
\begin{enumerate}
\item \textbf{Baseline (556–563s):} R fluctuates between 0.4–0.7.
\item \textbf{Coherence rise (563–566s):} R increases sharply from \textasciitilde{}0.5 to \textasciitilde{}0.85.
\item \textbf{Plateau (566–568s):} R stabilizes at \textasciitilde{}0.90–0.95—near-maximal global synchronization.
\item \textbf{Ignition (568s):} Brief fluctuation with peak power.
\item \textbf{Propagation (568–577s):} R remains elevated (0.70–0.90).
\item \textbf{Decay (577–582s):} Gradual decline to baseline.
\end{enumerate}
The plateau at R > 0.9 indicates that >90\% of electrode pairs are phase-synchronized—a highly organized state rarely sustained for multiple seconds.

\subsubsection{Information Flow Topology}

\textbf{Panel K: Directed Information Flow (Transfer Entropy Matrix).}
Asymmetric matrix showing z-scored change in transfer entropy (ΔTE) between all electrode pairs during ignition relative to baseline. Frontal sources (AF3, AF4, F3, F4) show strong positive ΔTE (z > 4) to most targets, indicating these regions become information sources during ignition. F8 shows particularly strong outward connections (z > 6), suggesting right frontal dominance in information broadcast. Parietal (P7, P8) and occipital (O1, O2) regions primarily receive information, consistent with frontal-to-posterior information flow.

\textbf{Panel L: Cross-Scale Information Flow (Δcorr² θ→γ).}
Time series showing the change in squared correlation between theta-band (4–8 Hz) and gamma-band (30–45 Hz) activity. A biphasic pattern emerges: pre-ignition surge (\textasciitilde{}+0.35 at 563–566s) indicates sudden theta-gamma coupling, followed by ignition-phase reversal (\textasciitilde{}−0.15 at 566–570s) suggesting decorrelation during maximal activity. This pattern suggests theta-gamma coupling initially increases to establish cross-frequency coordination, then partially decouples to allow independent processing before re-coupling during propagation.

\subsubsection{Exemplar Event Summary}
This high-quality SIE event demonstrates the complete set of signatures characteristic of the phenomenon:
\begin{table}[H]
    \centering
    \resizebox{\textwidth}{!}{%
    \begin{tabular}{lll}
        \toprule
        Feature & Observation & Implication \\
        \midrule
        \textbf{Multi-band activation} & All 6 $\varphi^n$ harmonics show $z > 2\sigma$ & Not theta-only phenomenon \\
        \textbf{Coherence-first} & PLV peaks 2--3s before z-scores & Network coordination precedes amplitude \\
        \textbf{Fundamental dominance} & $\Delta$HSI minimum precedes ignition & SR1-led organization established first \\
        \textbf{Bicoherence surge} & Multiple triads reach $>0.8$ & Nonlinear harmonic coupling \\
        \textbf{$\theta \rightarrow \gamma$ PAC increase} & MVL triples during ignition & Cross-frequency phase coupling \\
        \textbf{Global synchronization} & Kuramoto $R > 0.9$ at plateau & Whole-brain phase alignment \\
        \textbf{Frontal hub topology} & F8, AF3/4 show elevated $\Delta$TE & Directed information broadcast \\
        \textbf{State-space separation} & Distinct ignition vs.\ baseline clusters & Qualitatively different neural state \\
        \bottomrule
    \end{tabular}%
} \end{table}

The temporal sequence—fundamental dominance → theta-gamma coupling → coherence rise → power elevation → propagation → coordinated decay—suggests a hierarchical organization process that must be established before ignition can occur.

Having established the characteristic signatures of individual events, we now quantify these properties across the full population of 1,366 detected SIEs.

\subsection{Statistical Validation}
The preceding characterization of SIE events raises the question of whether these patterns reflect genuine neural organization or statistical artifacts. We address this through multiple complementary validation approaches.

\subsubsection{Cross-Device Validation}
One-way ANOVA tested whether harmonic frequencies differed across the three EEG recording devices (Emotiv EPOC X, n = 1,144; Muse, n = 165; Emotiv Insight, n = 57). Device effects varied by harmonic (Table 9): SR1 and SR3 showed no significant differences (p > 0.25), while SR5 exhibited a significant device effect (F = 5.74, p = 0.003), with Muse devices detecting systematically higher frequencies (32.88 Hz) than Emotiv devices (32.53–32.54 Hz).
% Requires: \usepackage{booktabs}
\begin{table}[H]
    \centering
    \begin{tabular}{llllll}
        \toprule
        Harmonic & EPOC X (Hz) & Muse (Hz) & Insight (Hz) & $F$ & $p$ \\
        \midrule
        SR1 & $7.63 \pm 0.33$ & $7.66 \pm 0.32$ & $7.68 \pm 0.32$ & 1.26 & 0.28 \\
        SR3 & $19.98 \pm 0.57$ & $20.03 \pm 0.57$ & $19.88 \pm 0.52$ & 1.34 & 0.26 \\
        SR5 & $32.53 \pm 1.23$ & $32.88 \pm 1.08$ & $32.54 \pm 1.41$ & 5.74 & 0.003 \\
        \bottomrule
    \end{tabular} \end{table}
For SR1 and SR3, cross-device frequency convergence (means within 0.3\% of grand mean) supports biological rather than device-specific origins. The significant SR5 device effect likely reflects differential hardware filtering or spectral sensitivity at higher frequencies (~32 Hz), where consumer-grade devices show greater variability. Importantly, the core $\phi^n$ ratios remain preserved across devices: the fundamental (SR1) and first ratio-determining harmonic (SR3) show device-independent detection. This validation is particularly notable given substantial hardware differences: channel count ranged from 4 (Muse) to 14 (EPOC X), electrode types included dry (Muse), semi-dry polymer (Insight), and saline-soaked felt (EPOC X), and spatial coverage varied from frontal-temporal (Muse) to comprehensive scalp coverage (EPOC X).

\subsubsection{Context Independence}
One-way ANOVA across five cognitive contexts (meditation, flow, non-flow, gaming, visual) revealed no significant frequency differences for any primary harmonic (Table 10).
% Requires: \usepackage{booktabs}
\begin{table}[H]
    \centering
    \begin{tabular}{llllllll}
        \toprule
        Harmonic & Meditation & Flow & Non-Flow & Gaming & Visual & $F$ & $p$ \\
        \midrule
        SR1 (Hz) & 7.67 & 7.66 & 7.62 & 7.61 & 7.61 & 1.46 & 0.21 \\
        SR3 (Hz) & 19.98 & 19.99 & 20.01 & 19.98 & 19.95 & 0.22 & 0.93 \\
        SR5 (Hz) & 32.74 & 32.56 & 32.57 & 32.52 & 32.40 & 1.68 & 0.15 \\
        \bottomrule
    \end{tabular}
\end{table}
The φⁿ architecture is preserved across contexts spanning eyes-closed meditation, cognitively demanding flow states, gaming engagement, and passive visual perception. This cross-context stability supports the interpretation that SIE frequencies reflect fundamental neural organizational constraints rather than task-specific processing.

\subsubsection{Null Control Validation}
To evaluate whether observed φⁿ precision is spectrally exceptional, we compared SIE harmonic triplets against random frequency triplets using two complementary null models.

\textbf{Null Model A: Uniform Sampling.} We generated 10,000 random frequency triplets by sampling uniformly from physiologically plausible ranges (SR1: 5–10 Hz; SR3: 17–23 Hz; SR5: 28–38 Hz). This yielded:
\begin{itemize}
\item Observed mean φ-error: 4.37\%
\item Random triplet mean error: 11.5\% ± 4.1\%
\item Cohen's d: 1.71 (very large effect)
\item p-value: 0.074
\end{itemize}

The marginal p-value (0.074) reflects a limitation of this null model: the sampling ranges are centered on φⁿ frequencies by construction $(7.5\,\text{Hz} \approx \varphi^0,\ 20\,\text{Hz} \approx \varphi^2,\ 33\,\text{Hz} \approx \varphi^3)$, biasing the null toward producing decent φ-ratios.

\textbf{Null Model B: Peak-Based Sampling.} To address this limitation, we implemented a more rigorous test using actual FOOOF peaks detected across all EEG sessions (N = 244,955 peaks from 661 recordings). Peaks were filtered into SR bands (SR1: 6.5–9.5 Hz, n = 23,075; SR3: 18–22 Hz, n = 27,659; SR5: 30–36 Hz, n = 18,347), and 10,000 random triplets were generated by sampling one peak from each pool. This controls for the actual spectral structure of EEG data rather than assuming uniform distributions.

\textbf{Peak-Based Results:}
\begin{itemize}
\item Observed mean φ-error: 4.37\% ± 2.00\%
\item Random peak triplet mean error: 9.05\% ± 4.15\%
\item Cohen's d: 1.44 (large effect)
\item p-value: $<$ 0.0001 (permutation test)
\item p-value: $<$ 0.0001 (Mann-Whitney U)
\end{itemize}

The peak-based null demonstrates that SIE events show \textbf{significantly better φ-precision} than random triplets sampled from actual EEG spectral peaks. The highly significant result (p $<$ 0.0001) with a large effect size (d = 1.44) indicates that the SIE detection algorithm preferentially identifies frequency combinations exhibiting exceptional φⁿ organization.

% Requires: \usepackage{booktabs}
\begin{table}[H]
    \centering
    \begin{tabular}{llll}
        \toprule
        Null Model & p-value & Cohen's d & Interpretation \\
        \midrule
        Uniform sampling & 0.074 & 1.71 & Marginal (biased null) \\
        \textbf{Peak-based} & \textbf{$<$ 0.0001} & \textbf{1.44} & \textbf{Highly significant} \\
        \bottomrule
    \end{tabular}
    \caption{Comparison of null control approaches for φⁿ precision validation.}
\end{table}

\begin{figure}[H]
    \centering
    \includegraphics[width=1\linewidth]{nc7_peak_based_results.png}
    \caption{Peak-based null control results. (A) Histogram comparing φ-error distributions for SIE events (red) versus random triplets sampled from actual FOOOF peaks (gray). SIE events cluster at lower φ-error values. (B) Box plot comparison showing significantly lower φ-error for SIE events. (C) Cumulative distribution functions demonstrating separation between distributions. (D) Percentile comparison across the error distribution.}
    \label{fig:nc7_peak_based}
\end{figure}

\subsubsection{Subject-Level Reproducibility}
To verify that findings are not driven by a subset of participants, we computed subject-level φⁿ ratio precision (N = 91 subjects).

\textbf{Subject-Level Results:}
\begin{itemize}
\item Mean SR3/SR1 ratio: 2.6254 ± 0.066 (predicted: 2.6180; error: +0.28\%)
\item Mean SR5/SR3 ratio: 1.6270 ± 0.045 (predicted: 1.6180; error: +0.56\%)
\end{itemize}

Subject-level ratio errors remained below 1\% for both primary ratios, confirming that φⁿ precision is not an artifact of averaging across heterogeneous subjects but represents a consistent organizational principle present within individual participants.

\subsection{Synchronization Metrics}
Having validated the robustness and specificity of SIE detection, we now characterize the synchronization properties that define these events.

\subsubsection{Power Elevation}
Power spectral elevation at SR harmonics during detected events was quantified using z-scores relative to session baselines. The fundamental frequency SR1 showed strongest elevation (z = 4.76 ± 2.95), with systematic decrease at higher harmonics (Table 11).
% Requires: \usepackage{booktabs}
\begin{table}[H]
    \centering
    \begin{tabular}{llll}
        \toprule
        Harmonic & Mean $z$ & Median $z$ & SD \\
        \midrule
        SR1 & 4.76 & 3.88 & 2.95 \\
        SR1.5 & 3.95 & 3.18 & 3.08 \\
        SR2 & 3.91 & 3.16 & 3.01 \\
        SR2o & 3.89 & 3.11 & 3.04 \\
        SR2.5 & 3.85 & 3.09 & 2.91 \\
        SR3 & 3.73 & 2.99 & 2.83 \\
        SR4 & 3.72 & 3.07 & 2.80 \\
        SR5 & 3.64 & 2.99 & 2.90 \\
        SR6 & 3.70 & 3.04 & 2.66 \\
        \bottomrule
    \end{tabular} \end{table}
All harmonics showed substantial power elevation (all median z > 2.7), confirming multi-band spectral enhancement during SIE events.

\subsubsection{Coherence and Phase-Locking}
Magnitude-squared coherence (MSC) and phase-locking value (PLV) quantified inter-regional synchronization at each harmonic (Table 12).
% Requires: \usepackage{booktabs}
\begin{table}[H]
    \centering
    \begin{tabular}{lll}
        \toprule
        Harmonic & Mean MSC & Mean PLV \\
        \midrule
        SR1 & 0.552 & 0.689 \\
        SR1.5 & 0.532 & 0.672 \\
        SR2 & 0.538 & 0.681 \\
        SR2o & 0.542 & 0.686 \\
        SR2.5 & 0.540 & 0.688 \\
        SR3 & 0.524 & 0.684 \\
        SR4 & 0.521 & 0.684 \\
        SR5 & 0.501 & 0.663 \\
        SR6 & 0.478 & 0.653 \\
        \bottomrule
    \end{tabular} \end{table}
Both MSC and PLV showed gradients decreasing from SR1 to SR6, with SR1 exhibiting highest synchronization (MSC = 0.552, PLV = 0.689). Notably, 60.5\% of events showed MSC > 0.5 and 81.5\% showed PLV > 0.6 at the fundamental frequency, confirming robust whole-brain synchronization during SIE events consistent with communication through coherence (\cite{fries2015rhythms}).

\subsubsection{Cross-Frequency Coupling}
Correlation analysis of z-scores across harmonics revealed strong cross-frequency coupling (Table 13).
% Requires: \usepackage{booktabs}
\begin{table}[H]
    \centering
    \begin{tabular}{ll}
        \toprule
        Pair & Pearson $r$ \\
        \midrule
        SR1 vs SR3 & 0.74 \\
        SR3 vs SR5 & 0.77 \\
        SR5 vs SR6 & 0.83 \\
        SR1 vs SR6 & 0.53 \\
        \bottomrule
    \end{tabular} \end{table}
MSC values showed even stronger coupling at higher harmonics (SR5–SR6: r = 0.87), suggesting coordinated synchronization across the harmonic stack during SIE events. This cross-frequency coupling in amplitude/coherence metrics contrasts with the independence observed in frequency values, indicating that SIEs involve coordinated power enhancement across independently-tuned oscillatory generators.

\subsection{Event Quality Assessment}
With the robustness of SIE detection established, we now examine what distinguishes high-quality from low-quality events.

\subsubsection{SR-Score Distribution}
The SR-Score (computed from sr1, sr3, sr5 with fixed weights) showed the following distribution across 1,366 events:
\begin{itemize}
\item Mean: 2.101 ± 1.879
\item Median: 1.599
\item Range: 0.060–30.152
\item Quartiles: Q1 = 1.032, Q2 = 1.599, Q3 = 2.570
\end{itemize}

\subsubsection{Harmonic Stack Index (HSI)}
HSI, measuring the ratio of overtone power to fundamental power, showed a median of 0.54. This indicates that summed overtone power was approximately 54\% of fundamental power during detected events—consistent with fundamental-dominated ignition, where the SR1 frequency drives the synchronized state while higher harmonics contribute secondary power.

\subsubsection{Relationship Between Quality and φ-Precision}
Correlation analysis revealed no significant relationship between SR-Score and φ² error magnitude (r = +0.052, p = 0.10). High-quality and low-quality events showed comparable φ-precision (high SR-Score: 4.27\% error; low SR-Score: 4.17\% error; t = 0.58, p = 0.56), indicating that golden ratio organization is a fundamental property of SIE events rather than an artifact of event quality or signal-to-noise ratio.

\subsubsection{Quality Determinants}
The SR-Score showed strong correlations with individual synchronization measures (Table 20).

\textbf{Table 20. SR-Score Correlations}
% Requires: \usepackage{booktabs}
\begin{table}[H]
    \centering
    \begin{tabular}{lll}
        \toprule
        Metric & Pearson $r$ & $p$-value \\
        \midrule
        $z$-score (SR1) & $+0.697$ & $< 10^{-198}$ \\
        MSC (SR1) & $+0.656$ & $< 10^{-168}$ \\
        PLV (SR1) & $+0.451$ & $< 10^{-68}$ \\
        HSI (inverse) & $-0.271$ & $< 10^{-23}$ \\
        \bottomrule
    \end{tabular}  \end{table}
Power elevation at SR1 was the strongest determinant of event quality (r = 0.697), followed by coherence (r = 0.656) and phase-locking (r = 0.451). This hierarchy suggests that amplitude increases at the fundamental frequency, supported by network-wide phase synchronization, are the primary characteristics distinguishing high-quality SIE events.

\subsection{Temporal Dynamics and Network Organization}
Beyond the static characterization of SIE properties, we examined the temporal evolution of ignition events to understand the dynamics underlying their formation.

\subsubsection{Six-Phase Temporal Evolution}
Individual SIE events exhibited a stereotyped temporal evolution characterized by six distinct phases (Figure 3). Analysis of Kuramoto order parameter R(t) trajectories across all 1,366 events revealed a canonical ignition sequence:

\begin{enumerate}
    \item \textbf{Phase 1 (Baseline):} Low-to-moderate global synchronization ($R \approx 0.4$--$0.6$) with characteristic fluctuations reflecting ongoing neural activity.
    \item \textbf{Phase 2 (Coherence Rise):} Rapid increase in phase alignment preceding amplitude changes, with R typically rising from baseline to > 0.7 over 1–3 seconds. This coherence-first signature is consistent with transient large-scale integration (\cite{varela2001brainweb}) rather than local oscillatory bursts.
    \item \textbf{Phase 3 (Plateau):} Brief stabilization (0.5–2 seconds) at elevated coherence levels (R > 0.8) immediately prior to amplitude ignition.
    \item \textbf{Phase 4 (Ignition):} Coincident peak in both coherence and power, marked by simultaneous z-score elevation across multiple harmonics. The fundamental frequency SR1 showed the most pronounced elevation (mean z = 4.57, 97.8\% of events exceeding 3σ).
    \item \textbf{Phase 5 (Propagation):} Sustained high-coherence state (5–15 seconds) with gradual spread of synchronization across harmonic frequencies. Cross-frequency coupling metrics (bicoherence, PAC) showed maximal values during this phase.
    \item \textbf{Phase 6 (Decay):} Gradual return to baseline over 3–8 seconds, with coherence typically declining before power—the inverse of the ignition sequence.
\end{enumerate}

\subsection{Emergent Frequency Properties}
Beyond the basic characterization of SIE events, analysis of the measured harmonic frequencies revealed unexpected mathematical structure.

\subsubsection{Golden Ratio (φⁿ) Frequency Organization}
Upon examining the nine harmonic frequencies detected across events (Table 4), a striking pattern emerged: the measured frequencies align closely with a geometric series f(n) = f₀ × φⁿ where φ = 1.618034 (the golden ratio). Using an empirically-derived fundamental frequency f₀ = 7.6 Hz, measured frequencies matched theoretical predictions across integer, half-integer, and quarter-integer values of n.

\textbf{Harmonic number (n) assignment methodology:} The assignment of specific n values to detected peaks was \textit{not} a post-hoc optimization but followed a constrained procedure: (1) the fundamental SR1 was assigned n = 0 by definition, anchoring the scale; (2) integer n values (1, 2, 3) were assigned to peaks nearest the canonical Schumann Resonance modes, preserving the physical correspondence to electromagnetic cavity modes; (3) half-integer and quarter-integer values were introduced only where robust spectral peaks existed at intermediate frequencies that did not match integer predictions.

Four of nine harmonics showed prediction errors below 1\%, with mean absolute error across all harmonics of 1.42\% and median error of 1.18\%. The integer harmonics (φ⁰, φ¹, φ², φ³) showed systematically better fit compared to half-integer harmonics, though this difference did not reach statistical significance. The quarter-integer harmonic SR2o (φ\textasciicircum{}1.25) showed excellent precision (0.76\% error), supporting the presence of fine-grained φⁿ structure.

\textbf{Optimal Fundamental Frequency.}
To identify the f₀ value providing best overall fit, we computed mean absolute prediction error across all harmonics for candidate f₀ values spanning 7.0–8.0 Hz (Figure 2A). Optimization yielded optimal f₀ = 7.58 Hz (minimizing squared error), 7.54 Hz (minimizing absolute error), and 7.58 Hz (N-weighted squared error).

\textbf{Table 5. Fundamental Frequency Comparison}
% Requires: \usepackage{booktabs}
\begin{table}[H]
    \centering
    \begin{tabular}{llll}
        \toprule
        $f_0$ Value (Hz) & Source & Mean $|$Error$|$ & Max $|$Error$|$ \\
        \midrule
        7.494 & Theoretical ($c/2\pi r$) & 1.57\% & 4.60\% \\
        7.540 & Optimal (min absolute) & 1.38\% & 3.96\% \\
        7.580 & Optimal (min squared) & 1.39\% & 3.30\% \\
        \textbf{7.600} & \textbf{Empirical (rounded)} & \textbf{1.45\%} & \textbf{3.14\%} \\
        7.630 & Measured SR1 mean & 1.49\% & 3.13\% \\
        7.830 & Schumann Resonance fundamental & 3.28\% & 5.61\% \\
        \bottomrule
    \end{tabular}
\end{table}
The empirical f₀ = 7.6 Hz provides near-optimal fit (1.45\% mean error) while maintaining interpretability. Critically, 7.6 Hz also corresponds to the empirical center of measured Schumann Resonance fluctuations: long-term monitoring data from the Space Observing System facility in Tomsk, Russia shows the SR fundamental varies around 7.6 ± 0.6 Hz across hours, days, and seasons (\cite{nickolaenko2014schumann}). Thus f₀ = 7.6 Hz represents both the empirically optimal fit to our EEG data and the measured center of SR fluctuation range. The canonical 7.83 Hz represents the modal peak of SR power spectral density, whereas 7.6 Hz represents the center of the frequency fluctuation window; our data suggest neural frequencies may track the latter rather than the former.

\textbf{Golden Ratio Precision in Harmonic Ratios.}
The φⁿ architecture predicts specific ratios between harmonics: SR3/SR1 = φ², SR5/SR1 = φ³, SR5/SR3 = φ, and SR6/SR4 = φ. These ratios are independent of f₀ choice and provide the most stringent test of golden ratio organization (Table 6).
% Requires: \usepackage{booktabs}
\begin{table}[H]
    \centering
    \begin{tabular}{lllll}
        \toprule
        Ratio & Predicted & Measured & Error (\%) & 95\% CI \\
        \midrule
        SR3/SR1 $= \varphi^2$ & 2.6180 & $2.6223 \pm 0.134$ & $+0.16$ & [2.614, 2.630] \\
        SR5/SR1 $= \varphi^3$ & 4.2361 & $4.2766 \pm 0.249$ & $+0.96$ & [4.262, 4.291] \\
        SR5/SR3 $= \varphi$ & 1.6180 & $1.6309 \pm 0.078$ & $+0.79$ & [1.627, 1.635] \\
        SR6/SR4 $= \varphi$ & 1.6180 & $1.6094 \pm 0.100$ & $-0.53$ & [1.604, 1.615] \\
        \bottomrule
    \end{tabular}
\end{table}
\textbf{Mean absolute ratio error: 0.61\%}
All four ratios deviated less than 1\% from predicted φⁿ values, with mean absolute error of 0.61\%. Bootstrap analysis (N = 10,000 iterations) confirmed that the 95\% confidence interval for SR3/SR1 (2.614–2.630) includes the predicted φ² value (2.618). However, the CI for SR5/SR3 (1.627–1.635) excludes the exact φ value (1.618), indicating a small but consistent positive deviation of approximately 0.8\%.

\textbf{Note on error metrics:} This 0.61\% represents the \textit{population-level} ratio error, computed using mean frequencies across all events (e.g., mean SR3 / mean SR1). Individual event-level precision is lower: when each event's φ-error is computed separately and then averaged across all 1,366 events, the mean composite φ-error is 4.37\%. This $\sim$7-fold difference is expected from the Central Limit Theorem: averaging $\sim$1,300 independent measurements reduces standard error by a factor of $\sqrt{1300} \approx 36$, yielding much higher precision for population means than for individual events. Both metrics are valid but measure different quantities: population-level precision (0.61\%) reflects the central tendency of the φⁿ architecture, while event-level precision (4.37\%) captures within-event variability.

\subsubsection{The Independence-Convergence Paradox}
A critical question concerns whether the three primary harmonics (SR1, SR3, SR5) vary together proportionally or independently. If frequencies were proportionally coupled (i.e., $\mathrm{SR3} \approx k \times \mathrm{SR1}$), we would expect strong positive correlations between harmonic frequencies across events. Instead, pairwise Pearson correlations revealed complete independence (Table 7).
% Requires: \usepackage{booktabs}
\begin{table}[H]
    \centering
    \begin{tabular}{llll}
        \toprule
        Pair & Pearson $r$ & $p$-value & Interpretation \\
        \midrule
        SR1 vs SR3 & $+0.030$ & 0.33 & No correlation \\
        SR1 vs SR5 & $-0.002$ & 0.94 & No correlation \\
        SR3 vs SR5 & $-0.020$ & 0.47 & No correlation \\
        \bottomrule
    \end{tabular} \end{table}
All correlations were statistically non-significant (all p > 0.3), with r values indistinguishable from zero. This demonstrates that harmonic frequencies vary completely independently across events—an event with high SR1 frequency is no more likely to have high SR3 or SR5 frequencies than expected by chance.

\textbf{Compensatory Error Mechanism.}
Despite frequency independence, harmonic ratios maintain remarkable precision (< 1\% error). This paradox resolves through analysis of ratio errors, which revealed a significant compensatory correlation (Table 8).
% Requires: \usepackage{booktabs}
\begin{table}[H]
    \centering
    \begin{tabular}{lll}
        \toprule
        Error Pair & Pearson $r$ & $p$-value \\
        \midrule
        SR3/SR1 error vs SR5/SR3 error & $\mathbf{-0.314}$ & $\mathbf{< 10^{-24}}$ \\
        SR3/SR1 error vs SR5/SR1 error & $+0.620$ & $< 10^{-100}$ \\
        SR5/SR1 error vs SR5/SR3 error & $+0.550$ & $< 10^{-80}$ \\
        \bottomrule
    \end{tabular} \end{table}
The critical finding is the strong negative correlation (r = −0.314, p < 10⁻²⁴) between SR3/SR1 error and SR5/SR3 error: when SR3/SR1 deviates above φ², SR5/SR3 compensatorily deviates below φ, and vice versa. This compensatory mechanism maintains overall φⁿ precision despite independent frequency variation, suggesting intrinsic mathematical constraints rather than coupled oscillatory generators.

\textbf{Population-Level φⁿ Encoding.}
To test whether ratio precision emerges from event-level harmonic coordination or population-level frequency distributions, we performed a shuffled bootstrap analysis. SR1, SR3, and SR5 frequencies were independently permuted across events (N = 10,000 iterations), destroying within-event pairings while preserving marginal distributions.

\textbf{Shuffled Bootstrap Results:}
\begin{itemize}
\item Observed mean composite φ-error: 4.37\%
\item Shuffled null mean error: 4.36\% ± 0.05\%
\item p-value: 0.60 (not significant)
\item Effect size (Cohen's d): −0.27
\end{itemize}

The non-significant p-value indicates that shuffling harmonic frequencies does not degrade ratio precision, implying that φⁿ relationships are encoded in the marginal frequency distributions themselves rather than through event-level coordination. The brain independently constrains each harmonic's frequency distribution to centers that satisfy φⁿ relationships at the population level.

\subsubsection{Frequency Stability Hierarchy}
Coefficient of variation (CV) analysis across harmonics revealed a systematic pattern of frequency stability. Among the primary harmonics (SR1–SR5), SR5 (β/γ boundary; CV = 1.79\%) showed the lowest variability, followed by SR3 (β boundary; CV = 2.02\%). The fundamental frequency SR1 showed highest variability (CV = 3.18\%). This pattern of decreasing CV from SR1 to SR5 suggests that higher harmonics show greater frequency stability than the fundamental, contradicting simple bottom-up models where fundamental frequency precision should propagate to higher harmonics. The beta-range harmonics (SR3, SR5) may serve as more stable reference frequencies than the theta-range fundamental (SR1), consistent with the two-subsystem model: a theta-centered system with moderate stability, and a tightly coupled beta-gamma complex with higher precision.

\subsubsection{Beat Frequency Architecture}
Inter-harmonic beat frequencies revealed additional structure consistent with φⁿ architecture. Both SR3−SR1 and SR5−SR3 beat frequencies approximate 12.4 Hz—the φ¹ frequency (α/β boundary)—suggesting nested oscillatory structure where inter-harmonic differences themselves follow φⁿ relationships. Specifically, SR3 − SR1 measured 12.38 ± 0.48 Hz (predicted: 12.13 Hz) and SR5 − SR3 measured 12.49 ± 0.70 Hz (predicted: 12.13 Hz).

\section{Discussion}

\subsection{Summary of Principal Findings}

The present study provides the first systematic characterization of Schumann Ignition Events (SIEs)—transient, high-coherence neural states exhibiting precise frequency organization at golden ratio ($\phi^n$) scaled harmonics. Analysis of 1,366 events across 91 participants, five cognitive contexts, and three EEG recording devices revealed several principal findings:

\begin{enumerate}
    \item \textbf{Reproducible neural phenomenon:} SIEs are characterized by brief (median 20 seconds) episodes of network-wide synchronization, with 81.5\% of events showing phase-locking values exceeding 0.6 and 60.5\% showing magnitude squared coherence exceeding 0.5 at the fundamental frequency.

    \item \textbf{Golden ratio frequency organization:} Harmonic frequencies spanning 7–40 Hz are organized according to powers of the golden ratio ($\phi$ = 1.618), with measured frequencies deviating less than 1.5\% from theoretical predictions $f(n) = f_0 \times \phi^n$ for an empirically-derived fundamental $f_0$ = 7.6 Hz.

    \item \textbf{Independence-convergence paradox:} Individual harmonic frequencies vary completely independently across events (all pairwise correlations $|r| < 0.03$), yet their population-level ratios maintain remarkable precision ($<$1\% error from $\phi^n$ predictions; event-level precision $\sim$4.4\%). This paradox suggests that golden ratio relationships are encoded in the marginal distributions of each harmonic rather than through event-level coordination.

    \item \textbf{Stereotyped temporal dynamics:} SIEs exhibit a six-phase temporal evolution in which coherence elevation precedes amplitude increases by 2–3 seconds, suggesting network-level coordination rather than local oscillatory bursts.

    \item \textbf{Cross-context and cross-device stability:} The $\phi^n$ frequency architecture is preserved across meditation, cognitive flow, non-flow arithmetic, gaming, and passive visual perception tasks, and across consumer EEG devices with substantially different electrode configurations, indicating a fundamental organizational property rather than task-specific artifact. Higher harmonics (SR5) showed modest device-related frequency shifts.
\end{enumerate}

\subsection{Relationship to Transient Synchronization Phenomena}

The temporal dynamics of SIEs share features with previously described transient synchronization phenomena. The coherence-first signature—phase alignment preceding amplitude elevation—is consistent with the ``ignition'' dynamics described in Global Neuronal Workspace theory (\cite{dehaene2011experimental, mashour2020conscious}). In this framework, conscious access is characterized by a late, large-scale ignition event involving prefrontal-parietal networks with sudden gamma-band enhancement and long-distance phase synchronization (\cite{gaillard2009converging}).

Similarly, the six-phase temporal evolution parallels the ``metastable'' dynamics observed in large-scale brain networks (\cite{tognoli2014metastable, deco2017dynamics}). Metastability refers to a dynamical regime intermediate between complete synchronization and complete independence, where periods of high synchronization alternate with desynchronized states in a continuous flow of transient coordination patterns. SIEs may represent particularly robust instances of such metastable coordination.

However, SIEs differ from previously characterized transient synchronization in several respects. First, the frequency specificity is unprecedented: rather than broadband synchronization or synchronization within conventional frequency bands, SIEs involve coordinated activity at nine discrete frequencies spanning theta through gamma that satisfy precise mathematical relationships. Second, the duration of 10–30 seconds exceeds typical gamma bursts (100–500 ms) while remaining shorter than sustained task-related oscillatory changes. Third, the cross-frequency coupling architecture—with a theta-anchored subsystem (SR1) showing moderate coupling to higher harmonics and a tightly coupled beta-gamma complex (SR3–SR6)—suggests hierarchical organization not captured by single-frequency synchronization accounts.

The discovery that SR-frequency organization exists continuously at baseline, with SIEs representing transient amplification rather than de novo generation, is consistent with the intrinsic ignition framework proposing that spontaneous brain dynamics constantly generate transient coordination events (\cite{deco2017dynamics}). This ``continuous-plus-amplification'' model has implications for understanding how the brain rapidly recruits coordinated network activity without requiring complete reorganization of oscillatory structure.

\subsection{The Independence-Convergence Paradox: A Novel Organizational Principle}

The most theoretically significant finding is the independence-convergence paradox: harmonic frequencies (SR1, SR3, SR5) vary completely independently across events (all correlations indistinguishable from zero), yet their ratios maintain less than 1\% deviation from golden ratio predictions. This paradox resolves through analysis showing a compensatory error mechanism ($r = -0.314$ between SR3/SR1 error and SR5/SR3 error) and shuffled bootstrap validation demonstrating that $\phi^n$ precision is preserved when frequencies are randomly permuted across events.

These findings indicate that golden ratio relationships are encoded in the population-level marginal distributions of each harmonic frequency, not through event-level coordination between oscillatory generators. Each harmonic appears to be independently constrained to a distribution whose mean satisfies $\phi^n$ relationships relative to other harmonics. This organizational principle differs fundamentally from both simple proportional coupling (where harmonic frequencies would covary) and independent oscillators (where ratio precision would degrade with frequency independence).

What biological mechanism could produce such population-level coordination among independently varying frequencies? One possibility is that each frequency band is constrained by local biophysical properties—membrane time constants, network architecture, receptor kinetics—that independently converge on $\phi^n$ values. The GABA-A receptor decay time constant ($\sim$25 ms for perisomatic basket cell synapses) directly determines gamma oscillation frequency through Pyramidal-Interneuron Network Gamma (PING) mechanisms (\cite{bartos2007synaptic, wang2010gamma, brunel2003fast}). Parvalbumin-positive fast-spiking interneurons, which constitute the primary cellular substrate for gamma generation (\cite{cardin2009driving, sohal2009parvalbumin}), show intrinsic resonance at approximately 40 Hz. If these timescales have been evolutionarily tuned to produce compatible frequency relationships—perhaps to optimize cross-frequency information transfer—independent variation within each system could still satisfy population-level $\phi^n$ constraints.

Alternatively, developmental processes might establish golden ratio relationships during circuit formation, with mature oscillatory generators retaining the mathematical structure while acquiring independent variability. Evidence that individual alpha frequency is highly heritable ($h^2 = 0.81$; \cite{smit2006heritability, posthuma2001heritability}) and stable across the lifespan despite substantial within-session variability (\cite{grandy2013individual}) supports the notion of genetically determined frequency ``set points'' that could encode $\phi^n$ architecture at the population level.

\subsection{Golden Ratio Organization and EEG Frequency Band Structure}

The $\phi^n$ frequency architecture provides a principled basis for understanding the organization of canonical EEG frequency bands (\cite{klimesch1999eeg}). Traditional band definitions (delta 1–4 Hz, theta 4–8 Hz, alpha 8–13 Hz, beta 13–30 Hz, gamma $>$30 Hz) have been established empirically over decades of research, but the boundaries between bands have remained somewhat arbitrary, with different researchers using different cutoffs and no underlying principle explaining why these particular frequency ranges should be functionally distinct.

The present findings suggest that EEG band boundaries may correspond to integer powers of $\phi$ (boundary frequencies), while band centers correspond to half-integer powers (attractor frequencies). Specifically: the theta-alpha boundary ($\sim$8 Hz) aligns with $\phi^0$ (7.6 Hz); the alpha-beta boundary ($\sim$12 Hz) aligns with $\phi^1$ (12.3 Hz); the beta boundary ($\sim$20 Hz) aligns with $\phi^2$ (19.9 Hz); and the beta-gamma boundary ($\sim$32 Hz) aligns with $\phi^3$ (32.2 Hz). The canonical alpha ``peak'' ($\sim$10 Hz) corresponds to the $\phi^{0.5}$ attractor (9.7 Hz); the sensorimotor rhythm ($\sim$15 Hz) corresponds to $\phi^{1.5}$ (15.6 Hz); and the 40 Hz gamma attractor corresponds to $\phi^{3.5}$ (40.9 Hz).

This framework aligns with large-scale spectral parameterization data showing that FOOOF-detected EEG peaks cluster at half-integer $\phi^n$ positions and are depleted at integer $\phi^n$ positions (\cite{donoghue2020fooof}), consistent with the attractor-boundary distinction. The framework also aligns with (\citet{buzsaki2004neuronal}) observation that mean brain oscillation frequencies form a linear progression on a natural logarithmic scale—the golden ratio ($\phi \approx 1.618$) is close to Euler's number ($e \approx 2.718$), and both predict approximately geometric frequency spacing.

The functional implications of $\phi^n$ organization extend to cross-frequency coupling. Phase-amplitude coupling between theta and gamma, alpha and gamma, and delta-theta and alpha-beta creates nested oscillatory hierarchies that transfer information across spatial and temporal scales (\cite{canolty2010role, jensen2007crossfrequency}). If these coupled frequencies satisfy $\phi^n$ relationships, non-integer ratios would minimize harmonic interference while maintaining mathematical precision—potentially optimizing information transfer across the frequency hierarchy (\cite{bonnefond2017communication}).

Intriguingly, the $\phi^n$ architecture exhibits self-similar structure in beat frequencies. Inter-harmonic beat frequencies (SR3 $-$ SR1 = 12.38 $\pm$ 0.48 Hz; SR5 $-$ SR3 = 12.49 $\pm$ 0.70 Hz) both approximate 12.4 Hz—the $\phi^1$ frequency corresponding to the alpha-beta boundary. This pattern suggests nested oscillatory structure where not only the harmonic frequencies themselves, but also their differences, follow $\phi^n$ relationships. Such recursive mathematical organization could provide multiple channels for cross-frequency information transfer: direct phase-amplitude coupling between harmonics, and amplitude-amplitude interactions at beat frequencies. The convergence of beat frequencies on $\phi^1$ may also explain why the alpha-beta transition zone ($\sim$12 Hz) is functionally significant across diverse cognitive tasks—it represents not just a boundary frequency but also the primary beat frequency generated by theta-beta and beta-gamma interactions.

\subsection{Two Synchronization Subsystems}

Analysis of cross-frequency amplitude correlations revealed a distinctive architecture suggesting two partially dissociable synchronization subsystems. The fundamental frequency SR1 (7.6 Hz) showed moderate coupling to higher harmonics ($r = 0.32$–$0.54$), while the beta-gamma harmonics SR3–SR6 (20–40 Hz) formed a tightly coupled complex (all pairwise $r > 0.6$). Importantly, these correlations reflect coordinated amplitude/power changes across harmonics—not frequency correlations, which remain near zero (see Independence-Convergence Paradox above).

This architecture aligns with established distinctions between theta-based and beta/gamma-based coordination mechanisms. Theta oscillations are associated with long-range hippocampal-cortical communication (\cite{hyman2005medial, colgin2009frequency}), working memory maintenance, and temporal organization of sequential information (\cite{sauseng2009brain}). Beta oscillations, in contrast, are associated with maintaining the ``status quo'' (\cite{engel2010beta, spitzer2017beyond}), sensorimotor integration, and long-range cortical synchronization where conduction delays challenge gamma coherence. The relatively weak coupling between SR1 and SR6 ($r = 0.32$) despite strong coupling within each subsystem suggests that theta-anchored and beta-gamma coordination mechanisms may be recruited somewhat independently during SIE ignition, with full integration occurring only at the population level through the $\phi^n$ constraint.

Having characterized the frequency organization and network architecture of SIEs, we now address the terminology and its relationship to geophysical phenomena.

\subsection{Relationship to Schumann Resonance}

The term ``Schumann Ignition Events'' references the correspondence between observed neural frequencies and the Schumann Resonances—extremely low-frequency electromagnetic phenomena arising from the Earth-ionosphere cavity, with a fundamental mode at approximately 7.83 Hz and higher harmonics at approximately 14.3, 20.8, 27.3, and 33.8 Hz (\cite{nickolaenko2014schumann}). Several aspects of this correspondence warrant careful interpretation.

First, the relationship is one of frequency correspondence, not demonstrated coupling. The measured Schumann Resonance fundamental (7.83 Hz) provides the poorest fit to observed neural frequencies (3.28\% mean error), while the theoretical electromagnetic prediction based on Earth's circumference ($c/2\pi r = 7.494$ Hz) and the empirically-derived optimal $f_0$ (7.54–7.60 Hz) provide better fits (1.38–1.45\% error). This suggests that neural $\phi^n$ architecture is not precisely tuned to measured SR frequencies, but rather to values approximately 3\% below the SR fundamental.

Second, without concurrent magnetometer recording of actual Schumann Resonance field strength, we cannot assess whether SIE occurrence correlates with SR activity or whether any causal relationship exists between planetary electromagnetic resonances and neural oscillations. Prior studies reporting correlations between SR intensity and human physiological variables have been limited by methodological concerns (\cite{cherry2002schumann, saroka2016similar, price2016schumann}), and proposed mechanisms for SR-brain coupling remain speculative given the extremely weak field strengths ($\sim$picoTesla) of SR signals.

Third, an alternative interpretation is that the frequency correspondence reflects convergent constraints rather than direct coupling. Both the Earth-ionosphere cavity (electromagnetic resonance) and the brain (neural oscillations) may independently optimize to similar frequencies due to shared physical principles—characteristic length scales, propagation velocities, and resonance conditions. The SR fundamental frequency is determined by $c/2\pi r$ where $c$ is the speed of light and $r$ is Earth's radius; interestingly, the theoretical neural $f_0$ providing best $\phi^n$ fit (7.54 Hz) is close to this electromagnetic prediction.

We therefore emphasize that ``Schumann'' in our terminology refers to the frequency range correspondence and the historical observation that human EEG and SR occupy overlapping spectral bands, not to proven electromagnetic coupling. Whether this correspondence reflects evolutionary adaptation to the planetary electromagnetic environment, coincidental convergence of physical constraints, or artifact of selection bias in EEG band definitions remains an open question requiring concurrent SR-EEG recording for resolution.

Regardless of the electromagnetic question, the neural phenomena themselves warrant functional interpretation.

\subsection{Functional Significance of SIEs}

What functional role might SIEs serve? Several lines of evidence suggest involvement in states of heightened cognitive integration. First, SIE detection rates varied substantially across cognitive contexts, with cognitive flow/non-flow tasks showing the highest rate (9.7 events/session) followed by meditation (7.6/session), gaming ($\sim$1.1/session), and passive visual perception ($\sim$1.1/session). While recording duration differences may contribute to this pattern, the elevated rates during flow states and meditation—both characterized by absorbed attention and reduced mind-wandering—suggest a relationship between SIEs and focused cognitive engagement.

Second, the temporal dynamics of SIEs—coherence preceding amplitude, fundamental dominance preceding ignition, and frontal-to-posterior information flow during peak synchronization—are consistent with top-down coordination mechanisms associated with attentional control and conscious access. The observation that high-quality events show dramatically stronger fundamental dominance (HSI 75\% lower, indicating the fundamental SR1 frequency dominates over overtones) and higher coherence (82\% higher) than low-quality events suggests that SR1-led frequency organization is functionally meaningful rather than epiphenomenal.

Third, the multi-band nature of SIEs—involving coordinated activity across nine harmonics spanning theta through gamma—distinguishes them from single-frequency phenomena and suggests involvement in multi-scale information integration. The hierarchical frequency organization described in predictive coding frameworks (\cite{bastos2012canonical, friston2009predictive}), where gamma carries feedforward prediction errors and alpha/beta carry feedback predictions, may require precisely-tuned frequency relationships to function optimally. SIEs may represent moments when this hierarchical architecture achieves maximal coordination.

Fourth, nonlinear cross-frequency coupling metrics reveal coordinated harmonic interactions during SIE events. Bicoherence analysis—which quantifies quadratic phase coupling between frequency triplets—showed multiple harmonic triads simultaneously surging to values exceeding 0.8 during the coherence/plateau phases, compared to baseline levels of 0.2–0.4. This synchronized increase across different frequency combinations indicates that nonlinear harmonic generation involves the entire $\phi^n$ stack through a common generative mechanism rather than independent oscillator coupling. Notably, bicoherence peaks slightly before maximal power elevation, suggesting that nonlinear coupling may drive harmonic power amplification rather than emerge as a consequence of it. Similarly, theta-gamma phase-amplitude coupling (PAC), measured by mean vector length, approximately triples during ignition (from $\sim$0.05 to $\sim$0.20), indicating substantial enhancement of hierarchical frequency organization consistent with nested oscillatory coding (\cite{lisman2013thetagamma}).

Fifth, examination of state-space dynamics reveals that SIE events occupy a qualitatively distinct neural regime. When plotting instantaneous coherence (Kuramoto order parameter) against harmonic power, baseline activity clusters at low power with moderate coherence (R $\sim$0.5–0.7), while ignition trajectories extend to high power while maintaining or increasing coherence (R $\sim$0.7–0.95), forming a characteristic ``lasso'' pattern. The clear non-overlap between baseline and ignition clusters demonstrates that SIEs are not simply extreme values of a continuous distribution but represent a separate attractor state in neural dynamics.

Sixth, the finding that power elevation at SR1 is the strongest determinant of event quality (r = 0.697), followed closely by coherence (r = 0.656), suggests that amplitude increases at the fundamental frequency, supported by network-wide phase synchronization, are the primary characteristics distinguishing high-quality SIE events. The strong but secondary role of coherence aligns with communication-through-coherence frameworks (\cite{fries2015rhythms}), where phase alignment between oscillatory populations gates the information flow that amplified signals carry.

Finally, the Harmonic Stack Index (HSI) findings—where fundamental dominance (decreasing HSI) precedes ignition by $\sim$3 seconds—suggest that SR1-led frequency organization may be a prerequisite for, rather than consequence of, ignition. The brain appears to first establish fundamental dominance over overtones before power can increase, implying that the SR1 frequency must lead the harmonic stack to support the subsequent synchronization cascade.

However, the functional significance of SIEs remains speculative without direct manipulation or behavioral correlates. Future studies should examine whether SIE occurrence predicts task performance, whether SIE frequency precision correlates with cognitive measures, and whether entrainment at $\phi^n$ frequencies enhances cognitive function relative to entrainment at non-$\phi^n$ frequencies.

\subsection{Methodological Considerations}

Several methodological factors warrant consideration in interpreting these findings.

The use of consumer-grade EEG devices (Emotiv EPOC X, Muse, Insight) with limited channel counts (4–14 channels) and lower signal-to-noise ratio compared to research-grade systems raises questions about spatial resolution and artifact contamination. However, these devices have been validated for research applications (\cite{badcock2015validation, krigolson2017choosing}).

Several observations argue against device-specific artifacts. First, cross-device frequency convergence was excellent, with means within 0.3\% of the grand mean for SR1 and SR3. Second, phase randomization—which preserves spectral content while destroying phase relationships—dramatically reduced synchronization metrics, indicating that the coherence signatures depend on genuine phase relationships rather than spectral artifacts. Third, the consistency of findings across devices with different electrode technologies (dry, semi-dry polymer, saline-soaked felt) and spatial coverage (frontal-temporal to comprehensive scalp) suggests biological rather than technical origins.

However, SR5 frequencies showed a significant device effect (F = 5.74, p = 0.003), with Muse devices detecting systematically higher frequencies than Emotiv devices ($\sim$0.35 Hz difference). This may reflect differences in anti-aliasing filters or spectral sensitivity at higher frequencies. While this does not invalidate the $\phi^n$ architecture—which is defined by frequency ratios rather than absolute frequencies—it indicates that absolute frequency claims for higher harmonics should specify the recording device.

The discovery phase relied on longitudinal recordings from a single experienced meditation practitioner, raising concerns about generalizability and potential overfitting of detection parameters. However, the validation phase demonstrated comparable SIE characteristics in 58 additional participants performing cognitive tasks unrelated to meditation, with no significant frequency differences between datasets.

The seven-stage detection pipeline involves multiple parameter choices (z-threshold, window duration, merging criteria, quality thresholds) that could influence detected event characteristics. Sensitivity analyses varying detection thresholds would strengthen confidence in the robustness of findings, though the large effect sizes observed ($d > 1.4$ for key comparisons) suggest that results would survive reasonable parameter variations.

\subsection{Limitations}

Several limitations constrain interpretation of the present findings. First, without concurrent magnetometer recording of Schumann Resonance field strength, we cannot assess whether SIE occurrence correlates with actual SR activity or whether any causal relationship exists. Claims about brain-environment electromagnetic coupling remain entirely speculative.

Second, the consumer EEG devices used provide limited spatial resolution. While the transfer entropy analysis suggested frontal hub topology during ignition, this conclusion is constrained by sparse electrode coverage and potential volume conduction effects. High-density EEG or source localization approaches would provide stronger evidence regarding the spatial organization of SIE generators.

Third, the present study is descriptive rather than mechanistic. We document the existence and characteristics of SIEs but cannot explain why neural oscillations would organize according to golden ratio mathematics. The mechanism generating $\phi^n$ architecture—whether evolutionary adaptation, developmental constraint, or emergent property of coupled oscillators—remains unknown.

Fourth, we did not examine behavioral correlates of SIE occurrence within events. Whether SIEs are associated with specific cognitive states, whether they predict task performance, and whether they are accessible to awareness remain open questions.

Fifth, the single-subject discovery dataset, while subsequently validated in independent samples, limits confidence that detection parameters generalize optimally across populations. Individual differences in baseline oscillatory dynamics, age effects on peak frequencies, and clinical variation could all influence SIE detection and characteristics.

\subsection{Future Directions}

Several directions would advance understanding of SIE phenomena. First, concurrent recording of EEG and Schumann Resonance field strength using magnetometers would enable direct assessment of brain-environment coupling hypotheses. If SIE occurrence correlates with SR field intensity or phase, this would support active coupling; if not, the frequency correspondence may reflect convergent constraints rather than electromagnetic interaction.

Second, high-density EEG with source localization would clarify the spatial organization of SIE generators. The present findings suggest frontal hub topology, but confirmation requires better spatial resolution than consumer devices provide.

Third, entrainment studies using transcranial alternating current stimulation (tACS; \cite{herrmann2016eeg, pogosyan2009boosting}) or sensory flicker (\cite{galambos1981human}) at $\phi^n$ versus non-$\phi^n$ frequencies would test whether the brain preferentially entrains to golden ratio frequencies. Prior work showing robust aftereffects of alpha-frequency tACS (\cite{zaehle2010alpha, kasten2016sustaining}) suggests that frequency-specific entrainment can produce lasting neural effects. If $\phi^n$ frequencies produce stronger entrainment or superior cognitive effects, this would support functional significance of the architecture.

Fourth, examination of SIEs in clinical populations—particularly those with altered oscillatory dynamics such as schizophrenia (gamma abnormalities; \cite{uhlhaas2008role, kwon1999gamma}), Parkinson's disease (beta abnormalities; \cite{brown2007basal, little2013movement}), or ADHD (theta abnormalities; \cite{barry2003review})—could reveal whether $\phi^n$ organization is disrupted in pathological states and whether restoration correlates with clinical improvement.

Fifth, cross-species validation would test whether $\phi^n$ architecture is evolutionarily conserved. Preliminary evidence suggests similar frequency organization across mammalian species despite 17,000-fold variation in brain volume (\cite{buzsaki2013scaling}); remarkably, even cephalopods—separated from vertebrates by 550 million years of evolution—show spindle-like 12–18 Hz oscillations during sleep (\cite{medendorp2024octopus}), suggesting these frequency bands represent convergent evolutionary solutions. Systematic examination of $\phi^n$ relationships across species would clarify whether this reflects fundamental biophysical constraints or convergent evolutionary optimization.

\subsection{Conclusions}

This study makes three principal contributions to our understanding of neural oscillatory organization.

\textbf{First, we establish SIEs as a reproducible neural phenomenon.} Across 1,366 events, 91 participants, five cognitive contexts, and three consumer EEG devices, Schumann Ignition Events showed consistent spectral, temporal, and network characteristics. The convergence of detection across devices with different electrode technologies and spatial coverage (4–14 channels) argues for biological rather than technical origins. The coherence-first temporal signature, frontal hub topology during ignition, and distinct state-space trajectories suggest that SIEs represent a qualitatively distinct neural regime rather than extreme values of continuous oscillatory dynamics.

\textbf{Second, we identify the independence-convergence paradox as a novel organizational principle.} Harmonic frequencies vary completely independently across events (all pairwise $|r| < 0.03$), yet their population-level ratios maintain $<$1\% deviation from golden ratio predictions. This paradox cannot be explained by coupled oscillators (which would show frequency covariation) or by independent rhythms (which would degrade ratio precision). Instead, it implies that each oscillatory generator is independently constrained to a frequency distribution whose mean satisfies $\phi^n$ relationships—a form of distributed coordination that preserves mathematical structure without requiring direct coupling. The compensatory error mechanism ($r = -0.314$ between successive ratio errors) further suggests an intrinsic constraint maintaining $\phi^n$ architecture.

\textbf{Third, we propose a unified $\phi^n$ framework for EEG frequency organization.} The observed harmonics map systematically onto canonical EEG bands: integer $\phi^n$ values correspond to band boundaries (theta-alpha at $\phi^0$, alpha-beta at $\phi^1$, beta-gamma at $\phi^3$), while half-integer values correspond to band centers or ``attractors'' (alpha peak at $\phi^{0.5}$, gamma at $\phi^{3.5}$). The self-similar beat frequency structure—where inter-harmonic differences also converge on $\phi^n$ values—suggests recursive mathematical organization that may optimize cross-frequency information transfer.

The correspondence between SIE frequencies and Earth's Schumann Resonances remains an open question. The neural fundamental (7.6 Hz) lies closer to the theoretical electromagnetic prediction ($c/2\pi r = 7.49$ Hz) than to the measured SR peak (7.83 Hz), suggesting that any relationship may reflect convergent physical constraints rather than direct electromagnetic coupling. Resolution requires concurrent EEG and magnetometer recording.

What the present findings establish is that transient, high-coherence neural states exhibit exceptional mathematical precision in their frequency organization—precision that emerges at the population level despite complete independence at the event level. This distributed encoding of golden ratio relationships across independently varying oscillators represents a novel form of neural coordination warranting further mechanistic investigation.

\printbibliography
\end{document}